\documentclass[a4paper, 10pt]{article}

\usepackage{packageList}
\usepackage{macros}

%

\graphicspath{{./Figures/}}
\pgfplotsset{compat=1.11}
\newlength\fwidth

\title{Sharp inverse estimates for \\ radial basis function interpolation: \\ One-to-one correspondence between \\ smoothness and approximation rates}

\author[1]{Tizian Wenzel \thanks{tizian.wenzel@mathematik.uni-stuttgart.de}}
%
%
%
\affil[1]{Institute for Applied Analysis and Numerical Simulation, University of Stuttgart, Germany}
%


\DeclareMathOperator{\spn}{span}
\DeclareMathOperator{\vol}{vol}


\DeclareMathOperator{\VN}{\tilde{V}_{N}}
\DeclareMathOperator{\proj}{\Pi_{\tilde{V}_{N}}}
\newcommand{\Ht}{\mathcal{H}_{\vartheta}}
\newcommand{\HT}{\mathcal{H}_{\Theta}}
\newcommand{\Nt}{{\tilde{N}}}

\newtheorem{conj}{Conjecture}

%
\newcommand{\tw}[1]{\textcolor{magenta}{TW: #1}}
\newcommand{\gs}[1]{\textcolor{blue}{GS: #1}}

%
\usepackage{comment}
%

%
\newif\iflong			
%
\longfalse

\newif\ifappendix
%
\appendixfalse

\usepackage{comment}
%

\usepackage[normalem]{ulem}			%


\makeatletter
\newcommand*\bigcdot{\mathpalette\bigcdot@{.5}}
\newcommand*\bigcdot@[2]{\mathbin{\vcenter{\hbox{\scalebox{#2}{$\m@th#1\bullet$}}}}}
\makeatother


\begin{document}

\maketitle %



  
\begin{abstract}
While direct statemets for kernel based approximation on domains $\Omega \subset \R^d$ using radial basis functions are well researched, 
there still seems to be a gap in comparison with corresponding inverse statements.
In this paper we sharpen inverse statements and by this close a ``gap of $d/2$" which was raised in the literature.
In particular we show that for a large class of finitely smooth radial basis function kernels such as Matérn or Wendland kernels, 
there exists a one-to-one correspondence between the smoothness of a function and its $L^2(\Omega)$ approximation rate via kernel interpolation: 
If a function can be approximation with a given rate, it has a corresponding smoothness and vice versa.
%
\end{abstract}

%
%
%
%
%
%
%
%
 

%
%
%
%
%
%
%
%

%
 
%
%
%
%
%
%
%
%
%
%
%
%
%
%
%
%
%
%
%
%
%
%
%
%
%
%
%
%
%
%
%
%
%
%
%
%
%
%
%
%
%
%


\section{Introduction} \label{sec:introduction}

In approximation theory, direct statements give error estimates for the approximation of an element (e.g.\ a function) %
using previously defined basis elements, e.g.\ polynomials.
The derived convergence rate usually depends on properties of the element to be approximated, e.g.\ the smoothness of the function.
Inverse statements treat the contrary and derive properties of the element which is approximated (e.g.\ the smoothness of the function) based on the approximation rate.

This paper deals with the particular case of function approximation via kernel interpolation, 
see e.g.~\cite{wendland2005scattered, fasshauer2007meshfree, fasshauer2015kernel, steinwart2008support} for a general introduction to kernel methods:
Given a nonempty set $\Omega$, a kernel is given as a symmetric function $k: \Omega \times \Omega \rightarrow \R$.
While the set $\Omega$ can be in general quite arbitrary, we are interested in the frequent case of domains $\Omega \subset \R^d$.
For a given set of $n \in \N$ points $X_n \subset \Omega$, the kernel matrix $A_{X_n}$ is defined as $(A_{X_n})_{ij} = (k(x_i, x_j))_{ij} \in \R^{n \times n}$.
If the kernel matrix is positive definite for any $n$ arbitrary pointwise different points $X_n \subset \Omega$ and any $n \in \N$, then the kernel is called strictly positive definite.
For such strictly positive definite kernels, there exists a unique associated Hilbert space $\ns$ of real-valued functions, the so called \textit{reproducing kernel Hilbert space} (RKHS), 
which is also called \textit{native space}.
In this space the kernel $k$ acts as a reproducing kernel, i.e.\ 
\begin{enumerate}
\item $k(\cdot, x) \in \ns ~ \forall x \in \Omega$,
\item $f(x) = \langle f, k(\cdot, x) \rangle_{\ns} ~ \forall x \in \Omega,  \forall f \in \ns$ \hfill (reproducing property).
\end{enumerate}
Kernels can be used for various tasks, and here we are concerned with the approximation of continuous functions by interpolation using continuous strictly positive definite kernels:
Given a continuous function $f \in \mathcal{C}(\Omega)$ and pairwise distinct interpolation points $X_n \subset \Omega$, we seek an interpolating function $s_{f, X_n}$ such that
\begin{align}
\label{eq:interpolation_conditions}
s_{f, X_n}(x_i) = f(x_i) \quad \forall x_i \in X_n.
\end{align}
A well known representer theorem shows
%
that there exists a minimum $\Vert \cdot \Vert_{\ns}$-norm interpolant of the form
\begin{align*}
s_{f, X_n}(\cdot) = \sum_{j=1}^n \alpha_j^{(n)} k(\cdot, x_j).
\end{align*}
The coefficients $\alpha_j^{(n)} \in \R, j= 1, \dots,n$ can be calculated by solving the linear equation system arising from Eq.~\eqref{eq:interpolation_conditions}, 
which is always solvable due to the strict positive definiteness of the kernel.

If the target function $f$ is contained in the RKHS of $k$, 
i.e.\ $f \in \ns$, then this interpolant $s_{f, X_n}$ coincides with the orthogonal projection of $f$ onto the subspace $V(X_n) := \Sp \{ k(\cdot, x_i),\; x_i \in X_n \}$, 
i.e.\ using the orthogonal projector ${\Pi_{V(X_n)}: \ns \rightarrow V(X_n)}$  it holds $s_{f, X_n} = \Pi_{V(X_n)}(f)$.

An important subclass of kernels is given by \textit{translational invariant kernels}, 
for which there exists a function $\Phi: \R^d \rightarrow \R$ such that it holds
\begin{align*}
k(x, z) = \Phi(x - z) ~ \text{for all} ~ x, z \in \R^d.
\end{align*}
A particular instance of this type of kernels is given by radial basis function kernels $k$ (RBF kernels), for which there exists a \textit{radial basis function} $\varphi: \R_{\geq 0} \rightarrow \R$ such that
\begin{align*}
k(x, z) = \varphi(\Vert x - z \Vert_2),
\end{align*}
which is again a translational invariant kernel with $\Phi(x) = \varphi(\Vert x \Vert_2)$.
In the following we will characterize the RBF kernel $k$ in terms of its Fourier transform $\hat{\Phi}$: 
We will assume the RBF kernel $k$ to be of finite smoothness $\tau > d/2$, 
which means there exists constants $c_\Phi, C_\Phi > 0$ such that 
\begin{align}
\label{eq:fourier_decay}
c_\Phi (1+\Vert \omega \Vert_2^2)^{-\tau} \leq \hat{\Phi}(\omega) \leq C_\Phi (1 + \Vert \omega \Vert_2^2)^{-\tau}.
\end{align}
In this case, the corresponding RKHS over the whole space $\R^d$ is norm equivalent to the Sobolev space of smoothness $\tau$, i.e.\ it holds $\mathcal{H}_k(\R^d) \asymp H^\tau(\R^d)$.
And under suitable conditions (such as a Lipschitz boundary) this norm equivalence extends also to domains $\Omega \subset \R^d$, i.e.\ it holds also $\ns \asymp H^\tau(\Omega)$.
The upper bound on the Fourier decay of Eq.~\eqref{eq:fourier_decay} usually yields direct statements, 
while the lower bound is important for inverse statements, see \Cref{subsec:kernel_approx}.
The kind of assumption from Eq.~\eqref{eq:fourier_decay} is frequently satisfied, e.g.\ for the class of Matérn or Wendland kernels.  %

In the literature, frequently the worst case error for kernel interpolation over all functions from the RKHS is analyzed.
For this, the \textit{power function} is considered, which is defined as
\begin{align}
\label{eq:def_powerfunc}
P_{X_n}(x) := \sup_{0 \neq f \in \ns} \frac{|(f - \Pi_{V(X_n)}(f))(x)|}{\Vert f \Vert_{\ns}}.
\end{align}
If the considered kernel satisfies Eq.~\eqref{eq:fourier_decay}, this approach yields $\Vert \cdot \Vert_{L^\infty(\Omega)}$ error estimates for $f \in \ns$ in the fill distance $h_{X, \Omega}$ (defined in Eq.~\eqref{eq:fill_sep_dist}) as
\begin{align}
\label{eq:Linfty_decay_error}
\Vert f - s_{f, X} \Vert_{L^\infty(\Omega)} \leq h_{X, \Omega}^{\tau - d/2} \cdot \Vert f - s_{f, X} \Vert_{\ns}.
\end{align}
%
%
%
%
%
%
%
%
%
In the following we will make use of the notation $s_{f, X}$ for the kernel interpolant, as we also want to interpolate functions $f \in \mathcal{C}(\Omega)$ which are not contained in the RKHS $\ns$, i.e.\ the projection property $s_{f, X} = \Pi_{V(X)}(f)$ does not hold. \\

One result on inverse estimates for kernel based interpolation was given in \cite{schaback2002inverse}, and we restate \cite[Theorem 6.1]{schaback2002inverse} here (with notation adapted to the current paper):
\begin{theorem}
\label{th:inverse_theorem_schaback}
Let $\Omega \subset \R^d$ be a bounded and open domain satisfying an interior cone condition.
The positive definite basis function $\Phi$ should satisfy the decay condition Eq.~\eqref{eq:fourier_decay}.
Suppose further that for some $f \in \mathcal{C}(\Omega)$ there exist constants $\mu > 0$ and $c_f > 0$ such that $\Vert f - s_{f, X} \Vert_{L^\infty(\Omega)} \leq c_f h_X^\mu$ for all $X \subset \Omega$ with $h_X$ sufficiently small.
If $\mu > \tau$, then $f$ must belong to the RKHS of $k$.
\end{theorem}
In comparison to the corresponding direct statement from Eq.~\eqref{eq:Linfty_decay_error},
\Cref{th:inverse_theorem_schaback} requires a decay rate larger than $\tau$ and thus leaves a ``gap of $d/2$", because the convergence rate required for the inverse statements is larger by $d/2$ than the convergence rate provided by the direct statement. 
One possible explanation for this ``gap of $d/2$" could be that $\Vert f - s_{f, X} \Vert_{\ns}$ provides a convergence rate of $h_{X, \Omega}^{d/2}$, see Eq.~\eqref{eq:Linfty_decay_error}.
However \cite[Section 6.2]{wenzel2022analysis} shows that this is in general not the case.

In the current work, we close this ``gap of $d/2$" by modifying \Cref{th:inverse_theorem_schaback}:
For this we consider the $L^2(\Omega)$ approximation error instead of the $L^\infty(\Omega)$ error and consider the worst case error over all possible sets on interpolation points, i.e.\ 
\begin{align*}
\sup_{\substack{X \subset \Omega \\ h_X \leq h}} \Vert f - s_{f, X} \Vert_{L^2(\Omega)}
\end{align*}
instead of the worst case error over all possible functions, as in Eq.~\eqref{eq:def_powerfunc}.
This change of objective for the worst case error enables us to derive precise error bounds for a fixed function $f \in \mathcal{C}(\Omega)$.
Furthermore our result extends to functions outside the RKHS (\textit{escaping the native space}), and by doing so we show a one-to-one correspondence between smoothness and approximation rate.
The main result on the inverse statements is given in \Cref{th:inverse_statement_generalized}, 
which gives rise to the one-to-one correspondence.


%


%
%
%
%
%
%
%
%
%
%
%
%
%

The paper is structured as follows: 
\Cref{sec:background} collects necessary background information on the tools which will be used. 
\Cref{sec:utility_statements} states and proves utility statements, 
which will then be leveraged in \Cref{sec:inverse_statement} to show inverse statements.
%
%
\Cref{sec:discussion_remarks} discusses the results and elaborates on open directions for future research,
while \Cref{sec:conclusion} concludes the paper.




%


%
%
%
%
%
%
%
%
%
%
%
%
%
%
%
%
%
%
%
%
%
%
%
%
%
%
%
%
%
%
%
%
%
%
%
%
%




\section{Background}
\label{sec:background}
%

%


\subsection{Geometric properties related to $\Omega$} \label{subsec:geometric_properties}

A standard assumption on the regularity of the domain, which is also frequently used in the following, is given by an interior cone condition.
We note that a Lipschitz boundary already implies the interior cone condition. 
%
\begin{definition}
A set $\Omega \subset \R^d$ is said to satisfy an interior cone condition, if there exists an angle $\alpha \in (0, \pi/2)$ and a radius $r>0$ such that for every $x \in \Omega$ a unit vector $\xi(x)$ exists such that the cone 
\begin{align}
\label{eq:definition_cone}
C(x, \xi(x), \alpha, r) := \{ x + \lambda y: y \in \R^d, \Vert y \Vert_2 = 1, y^\top \xi(x) \geq \cos(\alpha), \lambda \in [0, r] \}
\end{align}
is contained in $\Omega$.
\end{definition}
%

In order to quantify the distribution of scattered points $X \subset \Omega$ in a domain $\Omega \subset \R^d$, 
we recall the definition of separation distance $q_{X}$ and fill distance $h_{X, \Omega}$ as
\begin{align}
\label{eq:fill_sep_dist}
\begin{aligned}
q_{X} :=& \frac{1}{2} \min_{x_i \neq x_j \in X} \Vert x_i - x_j \Vert_2, \\
h_X := h_{X, \Omega} :=& \sup_{x \in \Omega} \min_{x_j \in X} \Vert x - x_j \Vert_2.
\end{aligned}
\end{align}
We call a sequence $(X_n)_{n \in \N}$ of sets of points $X_n \subset \Omega$ asymptotically uniformly distributed if the uniformity constant $\rho_{X_n} := \rho_{X_n, \Omega} := h_{X_n, \Omega} / q_{X_n}$ is uniformly bounded for all $n \in N$, i.e.\
\begin{align}
\label{eq:uniformly_distributed_points}
\exists C>0 ~ \forall n \in \N \quad \rho_{X_n} \equiv \frac{h_{X_n, \Omega}}{q_{X_n}} \leq C < \infty.
\end{align}
%
%
Also a sequence of points $( x_n )_{n \in \N}$ is called asymptotically uniformly distributed, if $( \{x_j\}_{j=1}^n )_{n \in \N}$ is asymptotically uniformly distributed.

Let $C_d$ denote the volume of the unit ball $B_1(0) \subset \R^d$ and let $C_{d, \alpha}$ be the volume of the unit cone $C(x, \xi(x), \alpha, 1)$ (which is independent of $x \in \Omega$), i.e.\
\begin{align}
\label{eq:constants_ball_cone}
\begin{aligned}
C_d &:= \vol(B_1(0)) \\
C_{d, \alpha} &:= \vol(C(x, \xi(x), \alpha, 1))
\end{aligned}
\end{align}
Then, based on geometric arguments for bounded domains $\Omega \subset \R^d$ which satisfy an interior cone condition, for a set $X_n \subset \Omega$ of $n$ points it holds
\begin{align}
\label{eq:geometric_bounds}
\begin{aligned}
h_{X_n, \Omega} &\geq c_{\Omega} \cdot n^{-1/d} \\
q_{X_n} &\leq C_{\Omega} \cdot n^{-1/d}
\end{aligned}
\end{align}
with constants
\begin{align*}
c_\Omega &= \frac{1}{\pi^{1/2}} \cdot \left( \vol(\Omega) \Gamma \left( \frac{d}{2} + 1 \right) \right)^{1/d} \\
C_\Omega &= \left( \frac{2\pi}{\alpha} \right)^{1/d} \cdot \frac{1}{\pi^{1/2}} \cdot \left( \vol(\Omega) \Gamma \left( \frac{d}{2} + 1 \right) \right)^{1/d}.
\end{align*}
(see \cite[Satz 2.1.6]{mueller2009komplexitaet} and \cite[Satz 2.1.7]{mueller2009komplexitaet}), especially it holds $C_\Omega > c_\Omega$. \\



A convenient way to obtain well distributed points for bounded domains $\Omega \subset \R^d$ is to use the \textit{geometric greedy algorithm}, which was introduced and analyzed in \cite{marchi2005optimal}.
We waive to recall the geometric greedy algorithm in itself, we simply recall that the resulting sequence of points satisfies for all $n \geq 2$ \cite[Lemma 5.1]{marchi2005optimal}
\begin{align}
\label{eq:geometric_uniformity}
\frac{1}{2} h_{X_n, \Omega} \leq \frac{1}{2} h_{X_{n-1}, \Omega} \leq q_{X_n} \leq h_{X_n, \Omega}.
\end{align}
%
%
The constructions within the next sections will frequently make use of these points:
As soon as a domain is fixed, we consider \textit{the} sequence of geometric greedy points, though it is clear that there can be several such sequences, e.g.\ depending on the choice of the first point.
However this will be used as notion to enforce that we think all the time about the same sequence of points.

In the following we use the shorthand notation $h_X := h_{X, \Omega}$, as the domain $\Omega$ is usually clear from the context.





\subsection{Kernel approximation}
\label{subsec:kernel_approx}

Additionally to the results presented in the introduction in \Cref{sec:introduction}, 
we recall direct statements for kernel interpolation yielding approximation rates as well as stability statements for kernel interpolation. 

We work under the assumption that $k$ is a RBF kernel that satisfies Eq.~\eqref{eq:fourier_decay} for some $\tau > d/2$ and that $\Omega$ has a Lipschitz boundary, 
i.e.\ the norm equivalence $\ns \asymp H^\tau(\Omega)$ holds.

We start with an direct statement, when the function $f \in \mathcal{C}(\Omega)$ of interest is not necessarily included in the RKHS: %

\begin{theorem}
\label{th:error_estimate}
Consider a RBF kernel $k$ that satisfies Eq.\ \eqref{eq:fourier_decay} for some $\tau > d/2$ and a bounded domain $\Omega \subset \R^d$ with Lipschitz boundary.
For some $\beta$ with $d/2 < \lfloor \beta \rfloor, \beta \leq \tau$
let $f \in H^\beta(\Omega) \supseteq H^\tau(\Omega)$.
Then it holds
\begin{align*}
\Vert f - s_{f,X} \Vert_{L^2(\Omega)} \leq C h_{X}^\beta \rho_{X}^\tau \Vert f \Vert_{H^\beta(\Omega)}.
\end{align*}
\end{theorem}

\begin{proof}
This is an immediate consequence of \cite[Theorem 4.2]{narcowich2006sobolev} using $\mu = 0$, 
with notation adopted to the current paper.
\end{proof}

\noindent For $\beta = \tau$ we have $f \in H^\tau(\Omega) \asymp \ns$, such that this is a classical RKHS result which can also be obtained by using the $\Vert \cdot \Vert_{L^\infty(\Omega)}$ decay of Eq.~\eqref{eq:Linfty_decay_error} and using Duchon's localization principle \cite{duchon1978erreur}, see e.g.\ \cite{wendland1997sobolev}.
However \Cref{th:error_estimate} goes further and also treats the case of \textit{escaping the native space}, i.e.\ $f \notin \ns$: 
It shows that 
-- when using well distributed interpolation points $X \subset \Omega$ such that $\rho_X$ is uniformly bounded -- 
the convergence rate $h_X^\beta$ for the approximation directly depends on the smoothness $\beta$ of the target function $f \in \mathcal{C}(\Omega)$.

Now we recall a stability statement related to the kernel matrix $A_X$:

\begin{theorem}
\label{th:stability_statement}
Let $k$ be a RBF kernel such that Eq.~\eqref{eq:fourier_decay} holds for some $\tau > d/2$. 
Then there is a $c_0 > 0$ such that for any set of pairwise distinct points $X \subset \Omega$ it holds
\begin{align*}
\lambda_{\min}(A_{X}) \geq c_0 q_{X}^{2\tau - d}
\end{align*}
and thus
\begin{align*}
\Vert A_X^{-1} \Vert_{2,2} \geq c_0^{-1} q_{X}^{d - 2\tau}.
\end{align*}
\end{theorem}

\begin{proof}
This is a well known result, see e.g.\ \cite{schaback1995error} or \cite[Chapter 12.2]{wendland2005scattered}.
\end{proof}

%
%
%
%
%
%
%
%
%
%
%
%
%
%
%
%
%
%
%


Additionally to the ``continuous" $L^2(\Omega)$ norms we consider the \textit{discrete} $L^2(X)$ norm for some $X \subset \Omega$, i.e.\ for $u \in \mathcal{C}(\Omega)$ %
%
\begin{align*}
%
\Vert u \Vert_{L^2(X)} \equiv \left( \frac{1}{|X|} \sum_{x \in X} |u(x)|^2 \right)^{1/2}.
\end{align*}

The following additional proposition allows to bound the ${\Vert \cdot \Vert_{\ns}}$ norm of an kernel interpolant via a matrix norm and a discrete $L^2(X)$ norm.
The proof is based on kernel interpolation properties and we include it here for convenience:

\begin{prop}
\label{prop:upper_bound_ns_norm_via_L2}
Let $X := \{x_1, ..., x_n\} \subset \Omega$. 
The $\Vert \cdot \Vert_{\ns}$ norm of $s_X := \sum_{j=1}^n \alpha_j k(\cdot, x_j)$ can be estimated via the discrete $L^2(X)$ norm:
\begin{align*}
\Vert s_X \Vert_{\ns}^2 &\leq \Vert A_X^{-1} \Vert_{2,2} \cdot |X| \cdot \Vert s_X \Vert_{L^2(X)}^2
\end{align*}
\end{prop}

\begin{proof}
We can calculate using $\alpha = (\alpha_1, ..., \alpha_n)^\top$:
\begin{align*}
\Vert s_X \Vert_{\ns}^2 &= \alpha^\top A_{X} \alpha = \alpha^\top A_{X} A_{X}^{-1} A_{X} \alpha 
\leq \Vert A_X^{-1} \Vert_{2,2} \cdot \Vert A_X \alpha \Vert_{2}^2 \\
& = \Vert A_X^{-1} \Vert_{2,2} \cdot |X| \cdot \Vert \sum_{j=1}^n \alpha_j k(\cdot, x_j) \Vert_{L^2(X)}^2 = \Vert A_X^{-1} \Vert_{2,2} \cdot |X| \cdot \Vert s_X \Vert_{L^2(X)}^2
\end{align*}
\end{proof}




%
%
%
%
%
%
%


\subsection{Powers of RKHS and interpolation spaces} \label{subsec:power_spaces}

In order to state our inverse statement in \Cref{th:inverse_statement_generalized} also for functions outside the RKHS, 
we recall the notion of powers of an RKHS and interpolation spaces and relate them to Sobolev spaces of intermediate smoothness.
This section is based on \cite{steinwart2012mercer, lions2012non, wendland2005scattered}.

Consider a strictly positive definite RBF kernel $k$ and a bounded domain $\Omega \subset \R^d$.
Let $T: L^2(\Omega) \rightarrow L^2(\Omega)$ be the kernel integral operator given by
\begin{align*}
Tv(x) := \int_{\Omega} k(x, y) v(y) ~ \mathrm{d}y, \qquad v \in L^2(\Omega), ~ x \in \Omega.
\end{align*}
This operator is then a compact operator, 
and Mercer's theorem gives the existence of a countable set of positive eigenvalues $\lambda_1 \geq \lambda_2 \geq .. > 0$ and corresponding continuous eigenfunctions $\{ \varphi_n \}_{n \in \N} \in L^2(\Omega)$ such that $T \varphi_n = \lambda_n \varphi_n$.
The set $\{ \varphi_n \}_{n \in \N}$ is an orthonormal basis for $L^2(\Omega)$ and an orthogonal basis in $\ns$. 
It holds the absolutely and uniformly convergent representation
\begin{align*}
k(x,y) = \sum_{n=1}^\infty \lambda_n \varphi_n(x) \varphi_n(y).
\end{align*}
With these tools the powers of the RKHS $\ns$ can be defined, which will also be called \textit{power spaces}: 
For $\vartheta \in [0, \infty)$ we define 
\begin{align}
\label{eq:power_spaces}
\Ht(\Omega) := \{ f \in L^2(\Omega): \sum_{j=1}^\infty \frac{|\langle f, \varphi_j \rangle_{L^2(\Omega)}|^2}{\lambda_j^\vartheta}  < \infty \} =\left\{
\begin{array}{ll}
L^2(\Omega) & \vartheta = 0 \\
\ns & \vartheta = 1 \\
TL^2(\Omega) & \vartheta = 2 \\
\end{array}
\right.
\end{align}
with dot products given as
\begin{align*}
\langle f, g \rangle_{\Ht(\Omega)} = \sum_{j = 1}^\infty \frac{\langle f, \varphi_j \rangle_{L^2(\Omega)} \langle g, \varphi_j \rangle_{L^2(\Omega)}}{\lambda_j^\vartheta}.
\end{align*}
The power spaces $\Ht(\Omega)$ for $\vartheta \in (0, 1)$ are exactly the interpolation spaces between $L^2(\Omega)$ and $\ns$ and (for $\vartheta$ sufficiently large) RKHS as well.
Obviously these power spaces $\Ht(\Omega)$ for $\vartheta \in (0, \infty)$ are nested, 
i.e.\ it holds $\mathcal{H}_{\vartheta_1}(\Omega) \supset \mathcal{H}_{\vartheta_2}(\Omega)$ for $\vartheta_1 < \vartheta_2$.
We have the following corollary, which is a direct consequence from the interpolation inequality (see e.g.\ \cite[Section 2.5]{lions2012non}) or can also be proven elementary via Hölder inequality:

\begin{cor}
\label{cor:hoelder_interpol}
Consider $\vartheta_1 < \vartheta < \vartheta_2$ and $g \in \mathcal{H}_{\vartheta_2}(\Omega)$.
Then it holds 
\begin{align}
\begin{aligned}
\label{eq:corr_hoelder}
\Vert g \Vert_{\mathcal{H}_{\vartheta}(\Omega)} &\leq \Vert g \Vert_{\mathcal{H}_{\vartheta_1}(\Omega)}^{\frac{\vartheta_2 - \vartheta}{\vartheta_2 - \vartheta_1}} \cdot \Vert g \Vert_{\mathcal{H}_{\vartheta_2}(\Omega)}^{\frac{\vartheta - \vartheta_1}{\vartheta_2 - \vartheta_1}} \\
\text{and especially \quad} \Vert g \Vert_{\mathcal{H}_{\vartheta}(\Omega)} &\leq \Vert g \Vert_{L^2(\Omega)}^{1 - \vartheta} \cdot \Vert g \Vert_{\ns}^{\vartheta},
\end{aligned}
\end{align}
whereby the second line is the special case $\vartheta_1 = 0, \vartheta_2 = 1$.
\end{cor}


\begin{figure}[t]
\label{fig:visualization}
\setlength\fwidth{.4\textwidth}
\begin{center}
\begin{tikzpicture}[>=latex, thick]
\draw[->] (0,0) -- (11cm,0) node [right] {$\R$};

\draw (1.5,-4pt) -- (1.5,4pt) node[below=10pt]{0};
\draw (1.5,5pt) node[above, align=center, fill=lightgray, rounded corners, inner sep=2pt]{$L^2(\Omega)$};

\draw (4,-2pt) -- (4,2pt);
\draw (4, 5pt) node[above, align=center, fill=lightgray, rounded corners, inner sep=2pt]{$\mathcal{H}_\vartheta(\Omega)$} node[below=10pt]{$\vartheta$};

\draw (5.5,-2pt) -- (5.5,2pt);
\draw (5.5, 5pt) node[above, align=center, fill=lightgray, rounded corners, inner sep=2pt]{$\ns$} node[below=10pt]{1};
\draw (9.5,-2pt) -- (9.5,2pt);
\draw (9.5, 5pt) node[above, align=center, fill=lightgray, rounded corners, inner sep=2pt]{$TL^2(\Omega)$} node[below=10pt]{2};
\end{tikzpicture}
\end{center}
\caption{Visualization of the scale of power spaces with the special cases $L^2(\Omega)$ (for $\vartheta = 0$), $\ns$ (for $\vartheta = 1$) and $TL^2(\Omega)$ (for $\vartheta = 2$).}
\end{figure}


%
%
Under suitable additional assumptions, the power spaces $\Ht(\Omega)$ for $\vartheta \in (0, 1)$ can be characterized in terms of Sobolev spaces:
Recall from \Cref{sec:introduction} that if the RBF kernel $k$ satisfies Eq.\ \eqref{eq:fourier_decay} and the domain $\Omega$ has a Lipschitz boundary, then it holds $\ns \asymp H^\tau(\Omega)$.
If the domain $\Omega \subset \R^d$ is additionally bounded, then the interpolation spaces between two Sobolev spaces $H^{\tau_1}(\Omega)$ and $H^{\tau_2}(\Omega)$ are the corresponding Sobolev spaces of intermediate smoothness, see \cite[Theorem 13.2.2]{agranovich2015sobolev} or \cite[Theorem 14.2.7]{brenner2008mathematical}.
As the norm equivalence $\ns \asymp H^\tau(\Omega)$ means that the identity operator $\mathrm{Id}_{\ns \rightarrow H^\tau(\Omega)}$ is a bounded operator, 
the norm equivalence $\ns \asymp H^\tau(\Omega)$ carries over to the corresponding interpolation spaces, i.e.\ we obtain for bounded Lipschitz domains $\Omega \subset \R^d$ and $\vartheta \in (0, 1)$ that $\Ht(\Omega) \asymp H^{\vartheta \tau}(\Omega)$.
The characterization of the power spaces $\Ht(\Omega)$ for $\vartheta > 1$ seems to be more challenging due to occuring boundary conditions in the space $TL^2(\Omega)$ \cite{schaback2018superconvergence}.
%


\subsection{Bernstein inequalities and related work} \label{subsec:related_work}

This subsection discusses related literature on inverse estimates and in particular Bernstein inequalities for kernel based interpolation:
The standard way to prove inverse statements for kernel-based approximation is to make use of Bernstein inequalities, 
which usually bound strong norms with help of a weaker norm times an inverse power of the separation distance.
A typical Bernstein inequality on a domain $\Omega \subseteq \R^d$ considers a kernel satisfying Eq.~\eqref{eq:fourier_decay} and a set of pairwise distinct points $X_n \subset \Omega \subset \R^d$ and is then e.g.\ of the form
\begin{align}
\label{eq:bernstein_not_available}
\left \Vert \sum_{j=1}^n \alpha_j \Phi(\cdot - x_j) \right \Vert_{H^s(\Omega)} \leq C q_{X_n}^{-s} \cdot \left \Vert \sum_{j=1}^n \alpha_j  \Phi(\cdot - x_j) \right \Vert_{L^2(\Omega)}
\end{align}
for some/all $s \in [0, \tau]$.
Such a Bernstein inequality Eq.~\eqref{eq:bernstein_not_available} would be helpful for the following and in particular for the proof of \Cref{th:inverse_statement_generalized},
however it is \textit{not} available in the literature in this general form to the best of the authors knowledge:

The statement \cite[Theorem 5.1]{narcowich2006sobolev} formulates and proves a Bernstein inequality,
however it considers $\Omega = \R^d$ instead of bounded domains $\Omega \subset \R^d$ and thus cannot be used to show inverse statements on $\Omega \subset \R^d$.
In \cite{ward2012lp}, these $\R^d$ Bernstein inequalities are extended to $L^p$ norms.
%
The statements in \cite[Section 6]{rieger2008sampling} generalize the $\R^d$ case to domains star-shaped domains $\Omega \subset \R^d$, 
however only for $s = \tau$.
However they require restrictive assumptions on the boundary and either restrict the point distribution $X \subset \Omega$  or require $\sum_{j=1}^n \alpha_j \Phi(\cdot - x_j)$ to be the interpolant of a function $f \in H^\tau(\Omega)$ or 
(see \cite[Theorem 6.4.1]{rieger2008sampling} and \cite[Corollary 6.5.3]{rieger2008sampling}).
The result \cite[Lemma 3.2]{cheung2018convergence} states another inverse inequality,
however the assumptions on the involved norms are not appropriate for our purposes.

The paper \cite{hangelbroek2018inverse} focusses on the case of domains $\Omega$, noting that these are substantially more difficult due to the occuring boundary.
By considering localized Lagrange functions, which are obtained by augmenting the set of interpolation points with centers outside the domain, inverse statements can be derived.
Subsequently, \cite{hangelbroek2018direct} extends these techniques and results to manifolds.

%
The spherical case $\Omega = \mathbb{S}^{d-1} \subset \R^d$ is better researched, 
as there are additional tools like spherical basis functions are available.
\cite[Theorem 6.2]{narcowich2007direct} is the corresponding statement to our main result \Cref{th:inverse_statement_generalized} for the spherical case.
\cite{mhaskar2010bernstein} deals as well with this particular case of the sphere and generalizes the Bernstein inequalities to the $L^p$ case,
and \cite[Lemma 9.2]{mirzaei2018petrov} generalizes these statements further.

As elaborated in the introduction, 
\cite[Theorem 6.1]{schaback2002inverse} (see \Cref{th:inverse_theorem_schaback}) shows an inverse statement and remarks that their key inequality can be seen as a Bernstein inequality, 
where they upper bound the RKHS norm $\Vert \cdot \Vert_{\ns}$ by an $\Vert \cdot \Vert_{L^\infty(\Omega)}$ norm. \newline

All in all, the literature so far does not provide a suitable Bernstein inequality to show our main result \Cref{th:inverse_statement_generalized} directly.
Instead of deriving a full Bernstein inequality, we leverage the statements provided in \Cref{sec:utility_statements}.


%



%
%
%
%
%
%
%
%
%
%
%
%
%
%
%
%
%
%
%
%
%
%
%
%
%
%
%
%
%
%
%
%
%
%
%
%
%
%
%
%
%
%
%
%
%
%
%
%
%
%
%
%
%
%
%
%
%
%
%
%
%
%
%
%
%
%
%
%
%
%
%




%
%
%
%
%
%
%
%
%
%
%
%
%
%
%




\section{Utility statements: From continuous $L^2(\Omega)$ to discrete $L^2(X)$ estimates}
\label{sec:utility_statements}

%

In this section we show how to convert estimates on $\Vert g \Vert_{L^2(\Omega)}$ into estimates on $\Vert g \Vert_{L^2(Y_1)}$ for continuous bounded functions $g \in \mathcal{C}(\Omega)$ and some well distributed set $Y_1 \subset \Omega$.
The basic connection between these two norms is given by the fact, that the integral of the $\Vert \cdot \Vert_{L^2(\Omega)}$ norm can be approximated by a quadrature given sufficiently many well distributed points.

In \Cref{th:utility_theorem_new} this is formalized, 
and then \Cref{prop:disc_L2_from_cont_L2_new} applies \Cref{th:utility_theorem_new} iteratively to even convert convergence rates on $\Vert f - s_{f, X_n} \Vert_{L^2(\Omega)}$ into convergence rates on $\Vert f - s_{f, X_n} \Vert_{L^2(X_{n+1})}$.
%
%

The proofs are quite constructive and leverage the points obtained from the geometric greedy algorithm, which was recalled in \Cref{subsec:geometric_properties} and also resembles ideas from related geometric constructions as in
\cite{duchon1978erreur, rieger2008sampling, rieger2010sampling}.
As these theorems are only used as a technical tool, the proofs are deferred to \Cref{sec:proofs}.

\subsection{Continuous $L^2(\Omega)$ estimates yield discrete $L^2(X)$ estimates}


The following \Cref{th:utility_theorem_new} shows that estimates on $\Vert g \Vert_{L^2(\Omega)}$ for $g \in \mathcal{C}(\Omega)$ imply estimates on $\Vert g \Vert_{L^2(Y_1)}$ for some well distributed set of points $Y_1 \subset \Omega$.
For reasons that will become clear later on when applying these results in \Cref{sec:inverse_statement}, we aim for sets $Y_1 \subset \Omega$ which are not "too close" to a given reference set $Y_0 \subset \Omega$.
The reason for this is that also the set $Y_0 \cup Y_1 \subset \Omega$ needs to be well distributed (in the sense of having a sufficiently large separation distance $q_{Y_0 \cup Y_1}$). 

It is worth to note that for $Y_0 = \{ \}$, i.e.\ without the reference set $Y_0$, the statement of \Cref{th:utility_theorem_new} is related to the upper bound within Marcinkiewicz-type discretization theorems \cite[Eq.\ (1.1)]{temlyakov2018marcinkiewicz},
however here we are interested only in a single given function $g \in \mathcal{C}(\Omega)$ instead of a whole subspace of functions.

\begin{theorem}
\label{th:utility_theorem_new}
Consider a bounded domain $\Omega \subset \R^d$ which satisfies an interior cone condition with angle $\alpha$ and radius $r > 0$. \\
For any bounded function $g \in \mathcal{C}(\Omega)$, any set $Y_0 \subset \Omega$ and any $q \leq \min \left( \frac{2^{-1/d} c_\Omega}{2 C_\Omega} \cdot q_{Y_0}, \frac{2}{5}r \right)$, 
there exists a set of points $Y_1 \subset \Omega$ with
\begin{align*}
\frac{1}{3} q \leq q_{Y_1} \leq h_{Y_1} \leq \frac{22}{3} \cdot q,
\end{align*}
and
\begin{align*}
\frac{1}{6}q \leq q_{Y_0 \cup Y_1} 
\end{align*}
such that
\begin{align*}
\frac{1}{|Y_1|} \sum_{y \in Y_1} |g(y)|^2 &\leq \tilde{C}_{d,\alpha} \cdot \int_{\Omega} |g(y)|^2 ~ \mathrm{d}y \\
\Leftrightarrow \qquad \Vert g \Vert_{L^2(Y_1)} &\leq \sqrt{\tilde{C}_{d,\alpha}} \cdot \Vert g \Vert_{L^2(\Omega)},
\end{align*}
with $\tilde{C}_{d,\alpha} := 4 \cdot  \frac{16^d C_d C_\Omega^d}{c_\Omega^{2d} C_{d, \alpha}^2}$ using $C_d, C_{d, \alpha}$ from Eq.~\eqref{eq:constants_ball_cone}.
\end{theorem}

The proof can be found in \Cref{sec:proofs}.


%
%
%


%

%
%
%
%
%
%
%
%
%
%
%
%
%
%
%
%
%
%
%
%
%
%
%
%
%
%
%
%
%
%
%
%
%
%
%
%
%
%
%
%
%
%
%
%
%
%
%
%
%
%
%
%
%
%
%
%
%
%
%
%
%
%
%
%
%
%
%
%
%
%
%
%
%
%
%
%
%
%
%
%
%
%
%
%
%
%
%
%
%
%
%
%
%
%
%
%
%
%
%
%
%
%
%
%
%
%
%
%
%
%
%
%
%
%
%
%
%
%
%
%
%
%
%
%
%
%
%
%
%
%
%
%
%
%
%
%
%
%
%
%
%
%
%
%
%
%
%
%
%
%
%
%
%
%
%
%
%
%
%
%
%
%
%
%
%
%
%
%
%
%
%
%
%
%
%
%
%
%
%
%
%
%
%
%
%
%
%
%
%
%
%
%
%
%
%
%
%
%
%
%
%
%
%
%


%
%
%

%
%
%
%
%
%
%
%
%
%
%
%
%
%
%
%
%
%
%
%
%
%
%
%
%
%
%
%
%
%
%
%
%
%
%
%
%
%
%
%
%
%
%
%
%
%
%
%
%
%
%
%
%
%
%
%
%
%
%
%
%
%
%
%
%
%
%
%
%
%
%
%
%
%
%
%
%
%
%
%
%
%
%
%
%
%
%
%
%
%
%
%
%
%
%
%
%
%
%
%
%
%
%
%
%
%
%
%
%
%
%
%
%
%
%
%
%
%
%
%
%
%
%
%
%
%

%
%
%
%

\subsection{Continuous $L^2(\Omega)$ convergence rates yield discrete $L^2(X)$ convergence rates}

In the following we consider the interpolation of continuous functions via kernel interpolation and their corresponding convergence rates in terms of the fill distance of the interpolation points.
The corresponding convergence rates are usually given in the continuous $\Vert \cdot \Vert_{L^2(\Omega)}$ norm, see \Cref{th:error_estimate}.

In this subsection we aim at using \Cref{th:utility_theorem_new} for $g := f - s_{f, X_n}$ in order to obtain convergence rates in a discrete $\Vert f - s_{f, X_n} \Vert_{L^2(X_{n+1})}$ norm for some sequence $(X_n)_{n \in \N} \subset \Omega$ based on convergence rates in the continuous $\Vert f - s_{f, X_n} \Vert_{L^2(\Omega)}$ norm.

\begin{prop}
\label{prop:disc_L2_from_cont_L2_new}
%
Consider a bounded domain $\Omega \subset \R^d$ which satisfies an interior cone condition with angle $\alpha$ and radius $r>0$. \\
Consider a function $f \in \mathcal{C}(\Omega)$ and a continuous kernel $k$ such that for some $c_f, \mu, h_0 > 0$ it holds
\begin{align*}
\sup_{\substack{X \subset \Omega \\ h_X < h}} \Vert f - s_{f, X} \Vert_{L^2(\Omega)} \leq c_f h^{\mu}
\end{align*}
%
%
%
%
for all $h \leq h_0$. \\
Then there exist constants $a \in (0, 1)$ and $c_0 > 0$ and a nested sequence of set of points $X_0 \subset X_1 \subset ... \subset X_m \subset .. \subset \Omega$ with 
\begin{align*}
\frac{1}{3} c_0 a^{m+1} \leq q_{X_{m+1}} \leq h_{X_{m+1}} \leq \frac{22}{3} c_0 a^{m+1}
\end{align*}
%
%
%
%
such that
\begin{align*}
\Vert f - s_{f, X_m} \Vert_{L^2(X_{m+1})} \leq \sqrt{\tilde{C}_{d,\alpha}} c_f h_{X_m}^{\mu}
\end{align*}
for all $m > 0$, using $\tilde{C}_{d,\alpha}$ defined in \Cref{th:utility_theorem_new}.
\end{prop}

The proof can be found in \Cref{sec:proofs}.
As it can be seen from that proof, also variations of the statement of \Cref{prop:disc_L2_from_cont_L2_new} with respect to the point distributions are easily derivable. 
We chose this formulation, as it is exactly what will be required in the following. \\
%




%
%
%
%
%
%
%
%
%
%







%
%
%


%
%
%
%
%
%
%
%
%
%
%
%
%
%
%
%
%
%
%
%
%
%
%
%
%
%
%
%
%
%
%
%
%
%
%
%
%
%
%
%
%
%
%
%
%
%
%
%
%
%
%
%
%
%
%
%
%
%
%
%
%
%
%
%
%
%
%






\section{Inverse statement: Continuous $L^2(\Omega)$ convergence rates yield smoothness} \label{sec:inverse_statement}
 %

In the following we show that convergence rates on $\Vert f - s_{f, X} \Vert_{L^2(\Omega)}$ in the fill distance $h_X$ imply a smoothness of the considered function $f$ in terms of the power spaces $\Ht(\Omega)$.
The proof idea roughly follows the one of \cite[Theorem 6.1]{schaback2002inverse}, 
but makes use of convergence rates in the $\Vert \cdot \Vert_{L^2(\Omega)}$ norm instead of the $\Vert \cdot \Vert_{L^\infty(\Omega)}$ norm and combines it with the theory of interpolation/power spaces.

%
%
%
%

\subsection{Case $\vartheta \in (0, 1]$: Escaping the native space} \label{subsec:vartheta_case_1}

The following inverse statement from \Cref{th:inverse_statement_generalized} is the first main result of this paper:

\begin{theorem}
\label{th:inverse_statement_generalized}
Consider a RBF kernel $k$ that satisfies Eq.\ \eqref{eq:fourier_decay} for some $\tau > d/2$ and a bounded domain $\Omega \subset \R^d$ which satisfies an interior cone condition with angle $\alpha$ and radius $r>0$.
Consider $f \in \mathcal{C}(\Omega)$ such that for some $c_f, \vartheta, h_0 > 0$ it holds
\begin{align}
\label{eq:inverse_statement_assumption}
\sup_{\substack{X \subset \Omega \\ h_X < h}} \Vert f - s_{f, X} \Vert_{L^2(\Omega)} \leq c_f h^{\vartheta \cdot \tau}
\end{align}
for all $h \leq h_0$. \\
If $\vartheta > 1$, then $f \in \ns$. If $\vartheta \in (0, 1]$,
then $f \in \mathcal{H}_{\vartheta'}(\Omega)$ for all $\vartheta' \in (0, \vartheta)$.
\end{theorem}

%
%
%
%
%
%
%
%


\begin{proof}
We start by using \Cref{prop:disc_L2_from_cont_L2_new} for $\mu := \vartheta \cdot \tau$:
Like this we obtain constants $a \in (0, 1)$ and $c_0 > 0$ and a nested sequence of sets of points $(X_m)_{m \geq 0}$ such that
\begin{align}
\begin{aligned}
\label{eq:proof:properties_sequence}
\frac{1}{3} c_0 a^{m+1} \leq& ~ q_{X_{m+1}} \leq h_{X_{m+1}} \leq \frac{22}{3} c_0 a^{m+1}, \\
%
\Vert f - s_{f, X_m} \Vert_{L^2(X_{m+1})} \leq& ~ \sqrt{\tilde{C}_{d,\alpha}} c_f h_{X_m}^{\mu}
\end{aligned}
\end{align}
Now we roughly follow the proof idea of \cite[Theorem 6.1]{schaback2002inverse}, but consider a weaker norm. %
We start by applying \Cref{cor:hoelder_interpol} to $g := s_{f, X_{n+1}} - s_{f, X_n} \in \ns$ using $\vartheta' < \vartheta$ for $\vartheta \in (0, 1]$ or $\vartheta' = 1$ for $\vartheta > 1$:
\begin{align}
\begin{aligned}
\label{eq:application_hoelder}
&\Vert s_{f, X_{m+1}} - s_{f, X_{m}} \Vert_{\mathcal{H}_{\vartheta'}(\Omega)}^2 \\
\leq& \Vert s_{f, X_{m+1}} - s_{f, X_{m}} \Vert_{L^2(\Omega)}^{2 - 2\vartheta'} \cdot \Vert s_{f, X_{m+1}} - s_{f, X_{m}} \Vert_{\ns}^{2\vartheta'} \\
\leq &\Vert s_{f, X_{m+1}} - s_{f, X_{m}} \Vert_{L^2(\Omega)}^{2 - 2\vartheta'} \\ 
&\qquad \cdot \Vert A_{X_{m+1}}^{-1} \Vert_{2,2}^{\vartheta'} \cdot |X_{m+1}|^{\vartheta'} \cdot \Vert s_{f, X_{m+1}} - s_{f, X_{m}} \Vert_{L^2(X_{m+1})}^{2\vartheta'},
\end{aligned}
\end{align}
whereby we estimated the $\Vert s_{f, X_{m+1}} - s_{f, X_{m}} \Vert_{\ns}$ quantity by using \Cref{prop:upper_bound_ns_norm_via_L2}.
%
%
%

For these four factors we have the following estimates:
%
\begin{itemize}
\item Adding $0 = f - f$, leveraging the assumed bound on $\Vert f - s_{f, X} \Vert_{L^2(\Omega)}$ and then estimating the fill distances as in Eq.\ \eqref{eq:proof:properties_sequence} gives
\begin{align*}
&\Vert s_{f, X_{m+1}} - s_{f, X_{m}} \Vert_{L^2(\Omega)} \\
\leq \, &\Vert f - s_{f, X_{m+1}} \Vert_{L^2(\Omega)} + \Vert f - s_{f, X_{m}} \Vert_{L^2(\Omega)} \\
\leq \, &2c_f h_{X_m}^{\mu} = 2c_f \left( \frac{22c_0}{3} \right)^\mu a^{\mu m}.
\end{align*}
\item For the norm of the inverse matrix we use \Cref{th:stability_statement} and subsequently the bound from Eq.\ \eqref{eq:proof:properties_sequence} on the separation distance to obtain
\begin{align*}
\Vert A_{X_{m+1}}^{-1} \Vert_{2,2} \leq c_0^{-1} q_{X_m \cup X_{m+1}}^{d-2\tau} \leq c_0^{-1} \left(\frac{1}{3} c_0 a^{m+1} \right)^{d-2\tau} 
= \frac{3^{2\tau-d}}{c_0^{2\tau - d + 1} a^{2\tau-d}} a^{-m (2\tau - d)}.
%
\end{align*}
\item The number of points $|X_{m+1}|$ can be upper estimated via Eq.\ \eqref{eq:geometric_bounds} as 
\begin{align*}
|X_{m+1}| \leq C_\Omega^d q_{X_{m+1}}^{-d} \leq \frac{3^d C_\Omega^d}{c_0^d} a^{-d(m+1)} = \frac{3^d C_\Omega^d}{c_0^d a^d} a^{-dm}.
\end{align*}
\item Using $s_{f, X_{m+1}}|_{X_{m+1}} = f|_{X_{m+1}}$ and then leveraging the bound on $\Vert f - s_{f, X_m} \Vert_{L^2(X_{m+1})}$ from Eq.\ \eqref{eq:proof:properties_sequence} and subsequently estimating the fill distances via Eq.\ \eqref{eq:proof:properties_sequence} gives
\begin{align*}
&\Vert s_{f, X_{m+1}} - s_{f, X_{m}} \Vert_{L^2(X_{m+1})} = \Vert f - s_{f, X_{m}} \Vert_{L^2(X_{m+1})} \\
%
\leq \, &c_f \cdot \sqrt{\tilde{C}_{d, \alpha}} \cdot h_{X_m}^\mu \leq \, c_f \cdot \sqrt{\tilde{C}_{d, \alpha}} \cdot \left( \frac{22}{3} c_0 a^m \right)^\mu \\
\leq \, &c_f \cdot \sqrt{\tilde{C}_{d, \alpha}} \cdot \left( \frac{22c_0}{3} \right)^\mu a^{\mu m}.
\end{align*}
\end{itemize}

\noindent Putting everything together we obtain
\begin{align*}
&\Vert s_{f, X_{m+1}} - s_{f, X_{m}} \Vert_{\mathcal{H}_{\vartheta'}}^2 \\
\leq &\,\Vert s_{f, X_{m+1}} - s_{f, X_{m}} \Vert_{L^2(\Omega)}^{2 - 2\vartheta'} \cdot \Vert A_{X_{m+1}}^{-1} \Vert_{2,2}^{\vartheta'} \cdot |X_{m+1}|^{\vartheta'} \cdot  \Vert s_{f, X_{m+1}} - s_{f, X_{m}} \Vert_{L^2(X_{m+1})}^{2\vartheta'} \\
\leq &\,C \cdot a^{2\mu m (1-\vartheta')} \cdot a^{-m (2\tau - d) \vartheta'} \cdot a^{-md \vartheta'} \cdot a^{2\mu m \vartheta'} \\
= &\,C \cdot a^{-m(-2\tau\vartheta + 2\tau \vartheta')} \\
= &\,C \cdot a^{2\tau m(\vartheta - \vartheta')},
\end{align*}
with constant $C > 0$ given by the collection of all the other occuring constants. %
%
%
%
For $\vartheta' < \vartheta$ we obtain $2\tau (\vartheta - \vartheta') =: \delta > 0$.
Then we can show that $(s_{f, X_m})_{m \geq 1}$ is a Cauchy series, i.e. for $n > m \geq m_0$:
\begin{align}
\label{eq:cauchy_sequence_calculation}
\Vert s_{f, X_{n}} - s_{f, X_m} \Vert_{\mathcal{H}_{\vartheta'}(\Omega)} &= \left \Vert \sum_{\ell=m}^{n-1} s_{f, X_{\ell+1}} - s_{f, X_{\ell}} \right \Vert_{\mathcal{H}_{\vartheta'}(\Omega)} \notag \\
&\leq \sum_{\ell=m}^{n-1} \Vert s_{f, X_{\ell+1}} - s_{f, X_{\ell}} \Vert_{\mathcal{H}_{\vartheta'}(\Omega)} \notag \\
&\leq C^{1/2} \cdot \sum_{\ell=m}^{\infty} a^{\ell\delta/2} 
\leq C^{1/2} \cdot \frac{a^{\delta m_0/2}}{1 - \sqrt{a}^{\delta}} \stackrel{m_0 \rightarrow \infty}{\longrightarrow} 0.
\end{align}
As $\mathcal{H}_{\vartheta'}(\Omega)$ is a complete space and $(s_{f, X_m})_{m \geq m_0}$ is a Cauchy sequence in $\mathcal{H}_{\vartheta'}(\Omega)$, 
there exists a unique limit element $\tilde{f} \in \mathcal{H}_{\vartheta'}(\Omega)$. \\
It remains to show that $\tilde{f}$ and $f$ coincide on $\Omega$. 
Therefore we consider
\begin{align*}
\Vert f - \tilde{f} \Vert_{L^2(\Omega)} %
&\leq \Vert f - s_{f, X_n} \Vert_{L^2(\Omega)} + \Vert \tilde{f} - s_{f, X_n} \Vert_{L^2(\Omega)} \\
&\leq c_f h_{X_n}^{\mu} + C_{\mathcal{H}_{\vartheta'} \hookrightarrow L^2(\Omega)} \Vert \tilde{f} - s_{f, X_n} \Vert_{\mathcal{H}_{\vartheta'}} \stackrel{n \rightarrow \infty}{\longrightarrow} 0.
\end{align*}
Thus we have $f = \tilde{f}$ in $L^2(\Omega)$ an $f$ can be seen as a particular representative of $\tilde{f}$, 
i.e.\ we have proven that $f \in \mathcal{H}_{\vartheta'}(\Omega)$ for any $\vartheta' < \vartheta$. 
\end{proof}

\noindent In the following we collect some remarks on the theorem and the proof:
\begin{itemize}
\item First we recall from \Cref{subsec:power_spaces} that under the additional assumptions of a Lipschitz boundary for $\Omega$, 
the spaces $\mathcal{H}_{\vartheta'}(\Omega)$ are normequivalent to the Sobolev spaces $H^{\vartheta' \tau}(\Omega)$.
Thus the smoothness characterization via power spaces is exactly the classical Sobolev smoothness.
\item Second we remark that the inequality in Eq.~\eqref{eq:application_hoelder} essentially bounds a strong norm (the $\Vert \cdot \Vert_{\mathcal{H}_{\vartheta'}(\Omega)}$ norm) via weaker norms (the $\Vert \cdot \Vert_{L^2(\Omega)}$ norm and the $\Vert \cdot \Vert_{L^2(X_m \cup X_{m+1})}$ norm, which can be again bounded via the $\Vert \cdot \Vert_{L^2(\Omega)}$ norm).
Thus this can be seen as a Bernstein inequality on general domains $\Omega \subset \R^d$ satisfying an interior cone condition.
However it is not a Bernstein inequality in its full generality, 
as it holds only with respect to the derived sequence $(X_m)_{m \geq m_0} \subset \Omega$ of  sets of points.
%
%
%
%
%
%
%
%

%
%
%
%
%
%
%
%
%
%
%
%
%
%
%
%
%
%
\item Third we remark that the statement of \Cref{th:inverse_statement_generalized} can be refined by additionally taking into account logarithmic decay rates, 
i.e.\ within \Cref{th:inverse_statement_generalized} we could assume
\begin{align*}
\sup_{\substack{X \subset \Omega \\ h_X < h}} \Vert f - s_{f, X} \Vert_{L^2(\Omega)} \leq c_f h^{\vartheta \cdot \tau} \cdot |\log(h)|^{\nu} \equiv c_f h^{\mu} \cdot |\log(h)|^\nu
\end{align*}
for some $\nu \in \R$.
Or on the other hand we could refine the decay of the Fourier transform within Eq.~\eqref{eq:fourier_decay} by taking into account corresponding logarithmic decay rates. \\
Due to $h_{X_m} \stackrel{\bigcdot}{\leq} a^m$ in the proof of \Cref{th:inverse_statement_generalized} (whereby $\stackrel{\bigcdot}{\leq}$ denotes inequality up to a constant), 
the decay $h_{X_m}^{\mu}$ turned into quantities like $a^{\mu m}$.
And thus the parts $|\log(h_{X_m})|^\nu$ will turn into $m^\nu |\log(a)|^\nu \stackrel{\bigcdot}{=} m^\nu$.
Therefore when showing the Cauchy sequence property in Eq.~\eqref{eq:cauchy_sequence_calculation}, 
we would not obtain the pure geometric series $\sum_{\ell=m}^{\infty} a^{\ell\delta/2}$ anymore.
Instead, in particular for $\vartheta = \vartheta'$, we would obtain a harmonic sum like $\sum_{\ell=m}^{\infty} \ell^{\tilde{\delta} \nu}$ (for some constant $\tilde{\delta} > 0$) which still converges for $\nu < 0$ small enough.
We waived to include this more detailed but more complicated analysis here in order to keep a clear outline of the exposition.
%
%
%

%
%
\end{itemize}

%




%
%
%
%

Finally we point out that in view of \Cref{th:error_estimate} the inverse statement in \Cref{th:inverse_statement_generalized} is sharp for $\Omega \subset \R^d$ with Lipschitz boundary:
For $f \in H^\beta(\Omega)$ with $d/2 < \lfloor \beta \rfloor, \beta \leq \tau$, 
the direct statement \Cref{th:error_estimate} provides an $\Vert \cdot \Vert_{L^2(\Omega)}$ error decay for asymptotically uniformly distributed points $X \subset \Omega$ according to $h_{X}^\beta$.
On the other hand, the inverse statement \ref{th:inverse_statement_generalized} already provides $f \in H^\beta(\Omega)$ if a convergence rate according to $h_{X}^\mu$ for any $\mu > \beta$ is realized. 

This connection is the one-to-one correspondence between smoothness and approximation rate.
In particular, no ``gap of $d/2$" is left anymore in comparison to \Cref{th:inverse_theorem_schaback} and the subsequent discussion of \Cref{sec:introduction}.





\subsection{Case $\vartheta \in (1, \infty)$: Superconvergence} \label{subsec:vartheta_case_2}

The case $\vartheta > 1$ for inverse statements is more challenging and left for future research.
We briefly comment on some challenges:

First, the corresponding direct statements (i.e.\ a statement like \Cref{th:error_estimate} for $\beta > \tau$) are still largely missing.
The only available results in this direction analyzes the error decay within $TL^2(\Omega) \subset \ns$ \cite{schaback2018superconvergence}.
However this space is difficult to analyze due to the occuring boundary conditions, as it frequently holds $TL^2(\Omega) \subsetneq H^{2\tau}(\Omega)$.
Second, the kernel approximant $s_{f, X_m}$ is not necessarily included in smoother spaces than $\ns$, 
therefore expressions like $\Vert s_{f, X_{m+1}} - s_{f, X_{m}} \Vert_{\Ht(\Omega)}$ for $\vartheta > 1$ or $\Vert s_{f, X_{m+1}} - s_{f, X_{m}} \Vert_{\mathcal{H}^{\tau'}(\Omega)}$ for $\tau' > \tau$ might not be well defined 
--- i.e.\ the current proof strategy of \Cref{th:inverse_statement_generalized} cannot be easily extended.







%
%
%
%
%
%
%
%



%
%
%
%
%
%
%
%
%
%
%
%
%
%
%
%
%
%
%
%
%
%
%
%
%
%
%
%
%
%
%
%
%
%
%
%
%
%
%
%
%
%
%
%
%
%
%
%
%
%
%
%
%
%
%
%
%
%
%
%
%
%
%
%
%
%
%
%
%
%
%
%
%
%
%
%
%
%
%
%
%


\section{Discussion and numerical example} \label{sec:discussion_remarks}

The starting point for this work was the ``gap of $d/2$" raised in \cite{schaback2002inverse}, 
namely the difference in the convergence rate for the direct statement in \Cref{th:error_estimate} and the corresponding inverse statement of \Cref{th:inverse_theorem_schaback}.
This ``gap of $d/2$" was closed in \Cref{th:inverse_statement_generalized} by considering the $L^2(\Omega)$ error instead of the $L^\infty(\Omega)$ error.
It remains an future challenge to transfer the $L^2(\Omega)$ error analysis of the inverse statement from \Cref{th:inverse_statement_generalized} into $L^\infty(\Omega)$ or other $L^p(\Omega)$ error statements that are sharp as well. 
Hereby we expect the $L^\infty(\Omega)$ error rates to be slower by a rate of $h_{X}^{d/2}$, 
see Eq.~\eqref{eq:Linfty_decay_error} and \Cref{th:error_estimate} (for $\tau = \beta$).
A possible tool to connect these might be Duchon's localization principle \cite{duchon1978erreur, wendland1997sobolev}.

\Cref{th:inverse_statement_generalized} makes use of well distributed points due to considering the fill distance $h_{X}$.
This is a frequent assumption in kernel based approximation, and also useful in applications since greedy algorithms as the $P$-greedy algorithm are known to provide well distributed points \cite[Theorem 19 \& 20]{wenzel2021novel}.
It remains open whether and how inverse statements can be generalized to adaptively distributed points.

The key ingredients for the proof of \Cref{th:inverse_statement_generalized} have been the technical statement of \Cref{prop:disc_L2_from_cont_L2_new} as well as standard RKHS analysis tools including stability estimates.
These tools do not seem to rely to much on the considered domain $\Omega \subset \R^d$, such that an extension to sphere or more general manifolds seems possible: 
In the case of manifolds, the asymptotic of fill and separation distance for well distributed points still coincide 
-- in view of Eq.~\eqref{eq:geometric_bounds} they will however depend on the effective dimension of the manifold instead of the dimension of the ambient space.
%
Also weakening the assumption of a Lipschitz boundary or cone condition might be possible, 
as \cite{wenzel2022stability} showed that the $L^\infty(\Omega)$ error estimates do not depend significantly on the shape of the domain $\Omega \subset \R^d$ 
-- though it remains open whether this holds as well for the $L^2(\Omega)$ error.

%
%
%
%
%
%
%
%
%
%
%
%
%
%

We continue with a remark on the considered kernels, i.e.\ RBF kernels that satisfy Eq.~\eqref{eq:fourier_decay}:
The reason for the one-to-one correspondence between smoothness and approximation rates as elaborated after \Cref{th:inverse_statement_generalized} is that the lower and upper bound on the Fourier transform within Eq.~\eqref{eq:fourier_decay} coincide:
The upper bound on the Fourier decay usually yields the error estimates (via the fill distance), i.e.\ the direct statements,
while the lower bound on the Fourier decay yields stability estimates (via the separation distance) which are required for the inverse statements \cite{schaback2002inverse, schaback1995error}.
In case of a gap between the lower and upper bound, i.e.
\begin{align*}
c_\Phi (1+\Vert \omega \Vert_2^2)^{-\tau_1} \leq \hat{\Phi}(\omega) \leq C_\Phi (1 + \Vert \omega \Vert_2^2)^{-\tau_2}
\end{align*}
for $d/2 < \tau_1 < \tau_2$,
this can result in a gap between direct and inverse statements.
Such a gap, i.e.\ $\tau_1 < \tau_2$, can happen if the Fourier transform $\hat{\Phi}$ has a staircase behaviour.

The RBF kernels satisfying Eq.~\eqref{eq:fourier_decay} are kernels of finite smoothness, as their RKHS is norm equivalent to Sobolev spaces.
However there are also kernels of infinite smoothness, such as the Gaussian kernel $k(x, y) = \exp(- \Vert x - y \Vert_2^2)$.
A one-to-one correspondence as discussed after \Cref{th:inverse_statement_generalized}) might not be derivable here:
For example for the Gaussian kernel, the known error estimates give rates as $\exp(c \log(h_{X}) / h_{X})$ \cite[Theorem 11.22]{wendland2005scattered}.
However the corresponding stability estimates look like $\exp(-40.71d^2/q_X^2) q_X^{-d}$ \cite[Corollary 12.4]{wendland2005scattered}, 
i.e.\ there is a large gap even for well distributed points, i.e.\ $q_X \asymp h_{X}$.
A detailed investigation on this possible \textit{analytic gap} seems to be an interesting future research direction.


%
%

%
%
%
%

Finally we conclude with a numerical example related to the one-to-one correspondence,
i.e.\ \Cref{th:error_estimate} and \Cref{th:inverse_statement_generalized}:

\begin{example} %
We consider the reconstruction of continuous functions on $\Omega := [0, 1]$ using the Matérn kernel $k(x, y) = \exp(\Vert x - y \Vert_2)$ which is an RBF kernel with $\Phi(r) = \exp(-r)$ and corresponding Fourier transform $\hat{\Phi}(\omega) = \sqrt{\frac{2}{\pi}} \cdot \frac{1}{\omega^2 + 1}$.
Thus the RKHS $\ns$ is norm equivalent to the Sobolev space $H^1(\Omega)$.
We consider the functions
\begin{align*}
f_\omega: \Omega \rightarrow \R, x \mapsto x^\omega \cdot (1-x)^\omega
\end{align*}
for $\omega > 0$ and it can be shown that $f_\omega \in H^{\omega + 1/2}(\Omega)$ for $\omega \in (0, \infty) \setminus \{ 1 \}$.
%
For the numerical experiment, we use $\omega \in \{ 0.01, 0.41, 0.81, 1.51, 2.01, 3.01 \}$ and consider the $\Vert f_\omega - s_{f_\omega, X} \Vert_{L^2(\Omega)}$ error using up to $2^{13} + 1 = 8193$ equidistant interpolation points $X \subset \Omega$.

The results are visualized in \Cref{fig:matern_L2_error}:
For $\omega \in \{ 0.01, 0.41 \}$ we have exactly the situation described in \Cref{th:error_estimate} and \Cref{th:inverse_statement_generalized}: 
We expect an decay of $n^{-0.51}$ respectively $n^{-0.91}$, which is indicated in \Cref{fig:matern_L2_error} via the dashed black lines.
For $\omega \in \{ 0.81, 1.51, 2.01, 3.01 \}$ we have the situation of \Cref{subsec:vartheta_case_2} which is not yet solved.
However in our case the convergence rate still seems to scale according to the smoothness up to some saturation point,
which seems to be reached for $\omega > 1.5$, i.e.\ as soon as $f_\omega \in H^{2\tau}(\Omega)$ (which is twice the smoothness of $\ns$).


\begin{figure}[hbt!]
\centering
\setlength\fwidth{.65\textwidth}
%
%
% This file was created by matlab2tikz.
%
%The latest updates can be retrieved from
%  http://www.mathworks.com/matlabcentral/fileexchange/22022-matlab2tikz-matlab2tikz
%where you can also make suggestions and rate matlab2tikz.
%
\definecolor{mycolor1}{rgb}{0.00000,0.44700,0.74100}%
\definecolor{mycolor2}{rgb}{0.85000,0.32500,0.09800}%
\definecolor{mycolor3}{rgb}{0.92900,0.69400,0.12500}%
\definecolor{mycolor4}{rgb}{0.49400,0.18400,0.55600}%
\definecolor{mycolor5}{rgb}{0.46600,0.67400,0.18800}%
\definecolor{mycolor6}{rgb}{0.30100,0.74500,0.93300}%
%
\begin{tikzpicture}

\begin{axis}[%
width=0.951\fwidth,
height=0.75\fwidth,
at={(0\fwidth,0\fwidth)},
scale only axis,
xmode=log,
xmin=1,
xmax=10000,
xminorticks=true,
ymode=log,
ymin=1e-10,
ymax=1,
yminorticks=true,
axis background/.style={fill=white},
axis x line*=bottom,
axis y line*=left,
legend style={at={(0.03,0.03)}, anchor=south west, legend cell align=left, align=left, draw=white!15!black}
]
\addplot [color=mycolor1, draw=none, mark=x, mark options={solid, mycolor1}]
  table[row sep=crcr]{%
5	0.395208116536507\\
9	0.276705579940274\\
17	0.19415136830309\\
33	0.136290595904114\\
65	0.0956707213365992\\
129	0.0671369233406873\\
257	0.0470816560337362\\
513	0.0329716899423768\\
1025	0.0230251719316315\\
2049	0.0159860591789946\\
4097	0.0109653073097013\\
8193	0.00732814015820774\\
};
\addlegendentry{$\alpha\text{ = 0.01}$}

\addplot [color=mycolor2, draw=none, mark=x, mark options={solid, mycolor2}]
  table[row sep=crcr]{%
5	0.0979824699649567\\
9	0.0504100680372608\\
17	0.026449545403499\\
33	0.0139888430151048\\
65	0.00742342557473034\\
129	0.00394524462201798\\
257	0.00209814876908024\\
513	0.00111611606878183\\
1025	0.000593654584748168\\
2049	0.000315488719096608\\
4097	0.000167135749131757\\
8193	8.76127864401554e-05\\
};
\addlegendentry{$\alpha\text{ = 0.41}$}

\addplot [color=mycolor3, draw=none, mark=x, mark options={solid, mycolor3}]
  table[row sep=crcr]{%
5	0.0234591475777655\\
9	0.00763129555394905\\
17	0.00265145688648768\\
33	0.000976421660622592\\
65	0.000374436069213298\\
129	0.000147083972470277\\
257	5.8533205646599e-05\\
513	2.34501353178688e-05\\
1025	9.42630404263753e-06\\
2049	3.79504319556154e-06\\
4097	1.52823000630411e-06\\
8193	6.13679011361556e-07\\
};
\addlegendentry{$\alpha\text{ = 0.81}$}

\addplot [color=mycolor4, draw=none, mark=x, mark options={solid, mycolor4}]
  table[row sep=crcr]{%
5	0.00630406485446773\\
9	0.0020747616300967\\
17	0.000612755422605411\\
33	0.000172328069068252\\
65	4.72041585384276e-05\\
129	1.27209707519712e-05\\
257	3.39006552478615e-06\\
513	8.96043324476924e-07\\
1025	2.3533990094565e-07\\
2049	6.14977398423333e-08\\
4097	1.60033049237389e-08\\
8193	4.14928945645237e-09\\
};
\addlegendentry{$\alpha\text{ = 1.51}$}

\addplot [color=mycolor5, draw=none, mark=x, mark options={solid, mycolor5}]
  table[row sep=crcr]{%
5	0.00438264237729036\\
9	0.00124196003898202\\
17	0.000319142077182057\\
33	8.02944248105614e-05\\
65	2.01022293570784e-05\\
129	5.02698669157969e-06\\
257	1.25679348976708e-06\\
513	3.14196164061406e-07\\
1025	7.85482744695472e-08\\
2049	1.96369352509154e-08\\
4097	4.90918146696508e-09\\
8193	1.22715266631218e-09\\
};
\addlegendentry{$\alpha\text{ = 2.01}$}

\addplot [color=mycolor6, draw=none, mark=x, mark options={solid, mycolor6}]
  table[row sep=crcr]{%
5	0.0013571525451123\\
9	0.00033662828460462\\
17	8.55054291097692e-05\\
33	2.14927143768166e-05\\
65	5.38095654761065e-06\\
129	1.3457326310639e-06\\
257	3.36464099458435e-07\\
513	8.41179596156526e-08\\
1025	2.10296102486533e-08\\
2049	5.25740776996481e-09\\
4097	1.31434301379472e-09\\
8193	3.28548176731015e-10\\
};
\addlegendentry{$\alpha\text{ = 3.01}$}

\addplot [color=black, dashed]
  table[row sep=crcr]{%
5	0.440073581512372\\
9	0.32608912857639\\
17	0.235760503367948\\
33	0.168096197994122\\
65	0.118963624669285\\
129	0.083868571853785\\
257	0.0590111539292018\\
513	0.0414801333411469\\
1025	0.0291427699910183\\
2049	0.0204698000849606\\
4097	0.0143761421499004\\
8193	0.010095877980592\\
};
% \addlegendentry{$\text{n}^{\text{-0.51}}$}

\addplot [color=black, dashed]
  table[row sep=crcr]{%
5	0.0693519298695475\\
9	0.0406219316539872\\
17	0.022772606209553\\
33	0.0124529879950983\\
65	0.00672001279836047\\
129	0.00360150994497902\\
257	0.00192345536414495\\
513	0.00102544991274774\\
1025	0.000546213635107033\\
2049	0.000290815849928053\\
4097	0.00015480223531499\\
8193	8.23925920369111e-05\\
};
% \addlegendentry{$\text{n}^{\text{-0.91}}$}

\addplot [color=black, dashed]
  table[row sep=crcr]{%
5	0.0121436514727618\\
9	0.00562266634314653\\
17	0.00244406027947407\\
33	0.0010250539859952\\
65	0.000421777591779447\\
129	0.000171841272297062\\
257	6.96606637807383e-05\\
513	2.81673669578976e-05\\
1025	1.13750086697238e-05\\
2049	4.59071185726992e-06\\
4097	1.8521219862933e-06\\
8193	7.47118910599067e-07\\
};
% \addlegendentry{$\text{n}^{\text{-1.31}}$}

\addplot [color=black, dashed]
  table[row sep=crcr]{%
5	0.008\\
9	0.00246913580246914\\
17	0.00069204152249135\\
33	0.000183654729109275\\
65	4.73372781065089e-05\\
129	1.20185085030948e-05\\
257	3.02805492891641e-06\\
513	7.59967929353381e-07\\
1025	1.90362879238548e-07\\
2049	4.76371837754421e-08\\
4097	1.19151103199213e-08\\
8193	2.9795047762132e-09\\
};
% \addlegendentry{$\text{n}^{\text{-2.00}}$}

\end{axis}
\end{tikzpicture}%	
%
\caption{Visualization of the $L^2(\Omega)$ error decay in the ($y$-axis) for interpolation with the Matérn kernel $k(x, y) = \exp(-\Vert x - y \Vert)$ on $\Omega = [0,1]$ in the number of equidistant interpolation points ($x$-axis) for different $\omega$ parameters of the functions $x^\omega \cdot (1-x)^\omega$.
This class of functions scales linearly in the Sobolev spaces $H^\tau(\Omega)$: 
It holds $x^\omega \cdot (1-x)^\omega \in H^{\omega + 1/2}(\Omega)$.
The dashed black lines indicate the convergence rates $n^{-0.51}, n^{-0.91}, n^{-1.31}$ and $n^{-2}$.}
\label{fig:matern_L2_error}
\end{figure}

\end{example}

\section{Conclusion} \label{sec:conclusion}

The current paper analyzed and closed an existing ``gap of $d/2$" between direct and inverse statements for radial basis function kernel interpolation.
This was enabled by deriving discrete $L^2(X)$ error estimates based of continous $L^2(\Omega)$ error estimates.
As a result, a one-to-one correspondence between smoothness (of the function to approximate) and the corresponding approximation rate was concluded for the \textit{escaping the native space case}, i.e.\ where the considered function is less smooth than the RKHS of the considered kernel.
A numerical experiment illustrated the result.
The derived techniques and results are expected to yield further consequences, e.g.\ for kernel quadratrue, solving PDEs with kernels or in statistical learning theory. \\
The results can be extended into several directions:
First, investigations of other kernels (non RBF kernels, analytic kernels) and other $L^p(\Omega)$ norms seems meaningful.
Second, the results should be extended to the superconvergence case, i.e.\ the case where the considered function is smoother than the RKHS. \\


\textbf{Acknowledgements:}
The author thanks Bernard Haasdonk and Gabriele Santin for support, discussions and feedback
and acknowledges funding of the project by the Deutsche Forschungsgemeinschaft (DFG, German Research Foundation) under Germany's Excellence Strategy - EXC 2075 - 390740016 and funding by the BMBF under contract 05M20VSA as well as support from the Studienstiftung des deutschen Volkes (German national Academic Foundation).
Part of the work was inspired/motivated by a research stay of the author with Lorenzo Rosasco at MaLGa, DIBRIS, Universit`a di Genova.
 
 
 

\IfFileExists{/home/ians1/wenzeltn/references.bib}{
\bibliography{/home/ians1/wenzeltn/references}      %
}{
\bibliography{/home/wenzeltn/references}			%
}
\bibliographystyle{abbrv} 
 
 
\appendix 

\section{Further proofs} \label{sec:proofs}

\subsection{Proof of \Cref{th:utility_theorem_new}}

\begin{proof}[Proof of \Cref{th:utility_theorem_new}]
As a first step, we approximate the integral $\int_\Omega |g(y)| ~ \mathrm{d}y$ with a quadrature:
Let $(y_j)_{j \in \N}$ be the sequence of geometric greedy points in $\Omega \subset \R^d$, $Y_N := \{y_1, ..., y_N\}$. 
\begin{itemize}
\item We cover $\Omega$ with patches $\{ E_i^{(N)} \}_{i=1}^N$, whereby $E_j^{(N)} \subset \Omega$ consist of all the points that are closer to $y_j$ than to any other $y_i, i=1, ..., N, i \neq j$:
\begin{align*}
E_j^{(N)} := \{y \in \Omega ~:~ \Vert y - y_j \Vert_2 \leq \Vert y - y_i \Vert_2 \quad \forall i \neq j \}
\end{align*}
Like this it follows for any $N \in \N$ that $\bigcup_{i=1}^N E_i^{(N)} = \Omega$ and $\mu(E_i^{(N)} \cap E_j^{(N)}) = 0$ for all $i \neq j$.
Furthermore all the patches $E_i^{(N)}, i=1, ..., N$ are asymptotically of the same size:
On the one hand, using that for any $j=1, ..., N$, 
the patch $E_j^{(N)}$ is included in a ball of radius $h_{Y_N, \Omega}$ around $y_j$, it holds
\begin{align}
\label{eq:patch_size_upper_bound}
\mu(E_j^{(N)}) \leq C_d h_{Y_N, \Omega}^d \stackrel{\text{Eq.\ } \eqref{eq:geometric_uniformity}}{\leq} C_d 2^d q_{Y_N}^d \stackrel{\text{Eq.\ } \eqref{eq:geometric_bounds}}{\leq} 2^d C_d C_\Omega^d N^{-1}.
\end{align}
On the other we have for any $j=1, ..., N$, 
\begin{align}
\label{eq:patch_size_lower_bound}
\mu(E_j^{(N)}) \geq C_{d, \alpha} q_{Y_N}^d \stackrel{\text{Eq.\ } \eqref{eq:geometric_uniformity}}{\geq} 
C_{d, \alpha} 2^{-d} h_{Y_N, \Omega}^d 
\stackrel{\text{Eq.\ } \eqref{eq:geometric_bounds}}{\geq}  2^{-d} c_\Omega^{d} C_{d, \alpha} N^{-1},
\end{align}
as the cone $C(y_j, \xi(y_j), \alpha, q_{Y_N})$ is associated to $y_j$ and the next closest point to $y_j$ has at least a distance of $2q_{Y_N}$.
\item For $N \rightarrow \infty$ it holds $\max_{i=1, ..., N} \sup_{a, b \in \overline{E_i^{(N)}}}( \Vert a - b \Vert_2) \rightarrow 0$, 
therefore for any integrable function $g \in \mathcal{C}(\Omega)$ ($g \in \mathcal{C}(\Omega)$ bounded provides integrability) the definition of the Rieman integral yields
\begin{align*}
\int_{\Omega} |g(y)| ~ \mathrm{d}y = \lim_{N \rightarrow \infty} \sum_{i=1}^N |g(y_i)| \mu(E_i^{(N)}).
\end{align*}
%
%
%
In particular we have
\begin{alignat*}{2}
&\forall \varepsilon > 0 ~ \exists N_\varepsilon \in \N ~ \forall N \geq N_\epsilon && \quad \left| \int_{\Omega} |g| ~ \mathrm{d}y - \sum_{i=1}^N |g(y_i)| \mu(E_i^{(n)}) \right| < \varepsilon \\
\Rightarrow \quad &\exists \tilde{N} \in \N ~ \forall N \geq \tilde{N}&& \quad \sum_{i=1}^N |g(y_i)| \mu(E_i^{(n)}) < 2 \cdot \int_{\Omega} |g| ~ \mathrm{d}y.
\end{alignat*}
Using Eq.\ \eqref{eq:patch_size_lower_bound} we obtain finally
\begin{align}
\label{eq:desired_estimate}
\forall N > \tilde{N}& \quad \frac{c_\Omega^{d} C_{d, \alpha}}{2^{d}} \cdot \frac{1}{N} \cdot \sum_{i=1}^N |g(y_i)| < 2 \cdot \int_{\Omega} |g| ~ \mathrm{d}y \notag \\
\Leftrightarrow \forall N > \tilde{N}& \quad \frac{1}{N} \cdot \sum_{i=1}^N |g(y_i)| < \frac{2^{d+1}}{c_\Omega^{d} C_{d, \alpha}} \cdot \int_{\Omega} |g| ~ \mathrm{d}y.
\end{align}
\end{itemize}

The set $Y_N$ of points can be quite huge as $\tilde{N}$ is not further specified.
Thus as a next step we want to construct smaller sets of points that still satisfy an inequality as in Eq.\ \eqref{eq:desired_estimate}: \\
For $q < \frac{c_\Omega}{2 C_\Omega} \cdot q_{Y_0}$ we define $Y'_1 := \{ y_1, y_2, ..., y_{n^*}\}$ with $n^*$ the smallest integer such that $h_{\{y_1, ..., y_{n^* + 1}\}} < q$.
From this definition we conclude
\begin{align}
\label{eq:proof:geom_props_1}
\frac{1}{2} q \leq \frac{1}{2} h_{Y'_1, \Omega} \leq h_{Y'_1 \cup \{ y_{n^*+1} \}, \Omega} < q
\end{align}
whereby the middle inequality is due to Eq.\ \eqref{eq:geometric_uniformity}.
For the separation distance we have
\begin{align}
\label{eq:proof:geom_props_2}
\frac{1}{2} q \leq \frac{1}{2} h_{Y_1', \Omega} \leq q_{Y_1'} \leq h_{Y_1', \Omega} < 2q.
\end{align}
With this definition we define $Y''_1 \subset Y'_1$ as the subset of all the points which are not too close to $Y_0$, i.e.\
\begin{align}
\label{eq:definition_y''}
Y''_1 := Y'_1 \cap \left( \Omega \setminus \bigcup_{y \in Y_0} B_{\frac{1}{2} q}(y) \right),
\end{align}
whereby $B_r(y) := \{ z \in \R^d ~ | ~ \Vert z - y \Vert_2 < r \}$, i.e.\ open balls. \\
The set $Y''_1$ is well distributed:
\begin{itemize}
\item For the separation distance $q_{Y''_1}$ we simply have 
\begin{align*}
q_{Y''_1} \geq q_{Y'_1} \geq \frac{1}{2} q.
\end{align*}
\item For the fill distance $h_{Y''_1, \Omega}$ we calculate:
\begin{align*}
h_{Y''_1, \Omega} \equiv& \sup_{y \in \Omega} \min_{y'' \in Y''_1} \Vert y - y'' \Vert_2 \\
=& \sup_{y \in \Omega} \min_{y'' \in Y''_1} \min_{y' \in Y'_1} \Vert y - y' + y' - y'' \Vert_2 \\
\leq& \sup_{y \in \Omega} \min_{y'' \in Y''_1} \min_{y' \in Y'_1} \Vert y - y' \Vert_2 + \Vert y' - y'' \Vert_2.
\end{align*}
Now pick $\tilde{y}' \in Y'_1$ as the closest $y'$ to $y$, i.e.\ $\tilde{y}' := \tilde{y}'(y) := \argmin_{y' \in Y'_1} \Vert y - y' \Vert_2$.
Then we continue to estimate:
\begin{align*}
h_{Y''_1, \Omega} \leq& \sup_{y \in \Omega} \left( \min_{y'' \in Y''_1} \Vert y - \tilde{y}'(y) \Vert_2 + \Vert \tilde{y}'(y) - y'' \Vert_2 \right) \\
\leq& \sup_{y \in \Omega} \left( \Vert y - \tilde{y}'(y) \Vert_2 + \min_{y'' \in Y''_1}  \Vert \tilde{y}'(y) - y'' \Vert_2 \right) \\
\leq& \sup_{y \in \Omega} \Vert y - \tilde{y}'(y) \Vert_2 + \sup_{y \in \Omega} \min_{y'' \in Y''_1}  \Vert \tilde{y}'(y) - y'' \Vert_2 \\
\leq& h_{Y'_1, \Omega} + \sup_{y \in \Omega} \min_{y'' \in Y''_1}  \Vert \tilde{y}'(y) - y'' \Vert_2.
%
%
%
%
%
%
\end{align*}
If $\tilde{y}'(y) \in Y_1''$, then the second summand is obviously zero.
If $\tilde{y}'(y) \in Y_1' \setminus Y_1''$, then by Eq.\ \eqref{eq:definition_y''} there is a point $y_0(y) \in Y_0$ such that $\Vert \tilde{y}'(y) - y_0(y) \Vert_2 < \frac{1}{2} q$. 
Then we choose any point $\tilde{y} := \tilde{y}(y) \in \Omega$ with $\Vert y_0(y) - \tilde{y}(y) \Vert_2 = \frac{5}{2} q$.
Such a point exists as $\frac{5}{2} q \leq r$, i.e. the distance is smaller than the radius of the cone condition of $\Omega$.
Adding $0 = - y_0(y) + y_0(y) - \tilde{y}(y) + \tilde{y}(y)$ gives
\begin{align*}
\sup_{y \in \Omega} \min_{y'' \in Y''_1}  \Vert \tilde{y}'(y) - y'' \Vert_2 
&= \sup_{y \in \Omega} \min_{y'' \in Y''_1}  \Vert \tilde{y}'(y) - y_0(y) + y_0(y) - \tilde{y}(y) + \tilde{y}(y) - y'' \Vert_2 \\
&\leq \sup_{y \in \Omega} \min_{y'' \in Y''_1}  \Vert \tilde{y}'(y) - y_0(y) \Vert_2 + \Vert y_0(y) - \tilde{y}(y) \Vert_2+ \Vert \tilde{y}(y) - y'' \Vert_2 \\
&\leq \sup_{y \in \Omega} \Vert \tilde{y}'(y) - y_0(y) \Vert_2 + \sup_{y \in \Omega} \Vert y_0(y) - \tilde{y}(y) \Vert_2 + \sup_{y \in \Omega} \min_{y'' \in Y''_1} \Vert \tilde{y}(y) - y'' \Vert_2 \\
&\leq \frac{1}{2} q + \frac{5}{2} q + 2q = 5q.
\end{align*}
In the last step we used for the last term that the closest point $y^*(y) \in Y'_1$ to $\tilde{y}(y)$ from $Y_1'$ has a distance of at most $2q$ (due to $h_{Y'_1, \Omega} < 2q$) and that this point must actually belong to $Y_1''$ (because $\Vert y_0(y) - y^*(y) \Vert_2 = \Vert y_0(y) - \tilde{y}(y) + \tilde{y}(y) - y^*(y) \Vert_2 \geq \left| \Vert y_0(y) - \tilde{y}(y) \Vert - \Vert \tilde{y}(y) - y^*(y) \right| > |\frac{5}{2} q - 2q | = \frac{1}{2} q$). \\
Thus finally we obtain
\begin{align*}
h_{Y_1'', \Omega} \leq 5q + 2q = 7q.
\end{align*}
\end{itemize}
For ease of notation, we set $Y_1'' =: \{y''_1, ..., y''_{|Y_1''|} \}$.
Now we create sets by assigning nearby points from $Y_N$ (for $N$ large enough) to the points from $Y_1''$:
\begin{align}
\label{eq:proof:definition_TiN}
T_i^{(N)} :=& \{ z \in Y_N ~ | ~ \Vert y''_i - z \Vert_2 \leq \frac{1}{3} q \}, \qquad &&i=1, ..., |Y_1''|.
\end{align}
The distance $\frac{1}{3}q$ was chosen such that $T_i^{(N)} \cap T_j^{(N)} = \emptyset$ for any pairwise distinct $i, j = 1, ..., |Y_1''|$. 
For convenience we introduce
\begin{align*}
T_i^{(N)} =: \left\{ y_1^{(i, N)}, ..., y_{|T_i^{(N)}|}^{(i, N)} \right\}.
\end{align*}
Now we prove that for $N \rightarrow \infty$, all the sets $T_i^{(N)}$ are of asymptotically equal size:
\begin{itemize}
\item The number of points must be large enough, to at least allow to cover the surrounding of $y_i'' \in Y_1''$ in terms of the volume:
The surrounding of $y_i''$ consist at least of the associated cone $C(y_i'', \xi(y_i''), \alpha, \frac{1}{3} q)$:
%
\begin{align*}
\bigcup_{y \in T_j^{(N)}} B_{h_{Y_N, \Omega}}(y) &\supset C(y_i'', \xi(y_i''), \alpha, \frac{1}{3}q) \\
\Rightarrow \quad \#T_j^{(N)} \cdot C_d h_{Y_N, \Omega}^{d} &\geq \mu(C(y_j'', \xi(y_j''), \alpha, \frac{1}{3}q)) = 3^{-d} C_{d, \alpha} \cdot q^d.
\end{align*}
Rearranging and using the geometric bounds from Eq.\ \eqref{eq:geometric_bounds} and Eq.\ \eqref{eq:geometric_uniformity} gives
\begin{align}
\label{eq:subset_size_lower_bound}
\# T_i^{(N)} &\geq \frac{3^{-d} C_{d, \alpha}}{C_d} \cdot \frac{q^d}{h_{Y_N, \Omega}^d}
\geq \frac{3^{-d} C_{d, \alpha}}{C_d} \cdot \frac{q^d}{2^d C_\Omega^d N^{-1}} \notag \\
&= \frac{C_{d, \alpha} q^d}{6^d C_d C_\Omega^d} \cdot N.
\end{align}
\item Using again a volume comparison argument, we have (for $N > n$ large enough such that $q_{Y_N} \leq \frac{2}{3} q$)
\begin{align*}
%
\bigcup_{y \in T_j^{(N)}} B_{q_{Y_N}}(y) &\subset B_{\frac{1}{3}q + q_{Y_N}}(y''_j) \\
\Rightarrow \quad \#T_j^{(N)} \cdot C_d q_{Y_N}^d &\leq C_d \cdot \left( \frac{1}{3}q + q_{Y_N} \right)^d \leq C_d q^d
\end{align*}
Using the geometric bounds from Eq.\ \eqref{eq:geometric_bounds} and the asymptotical uniform distribution of Eq.\ \eqref{eq:geometric_uniformity} we arrive at
\begin{align}
\label{eq:subset_size_upper_bound}
\#T_j^{(N)} &\leq q^d \cdot q_{Y_N}^{-d} \leq q^d 2^d c_\Omega^d N.
\end{align}
\item Combining Eq.\ \eqref{eq:subset_size_lower_bound} and Eq.\ \eqref{eq:subset_size_upper_bound} formalizes our claim that the sets $T_j^{(N)}$ are of asymptotical equal size, 
because the lower and upper bound do actually not depend on $j = 1, ..., |Y_1''|$:
\begin{align}
\label{eq:subset_size_asymp}
\frac{C_{d, \alpha}}{6^d C_d C_\Omega^d} \cdot q^d N \leq 
\#T_j^{(N)} 
\leq 2^d c_\Omega^d \cdot q^d N.
\end{align}
\end{itemize}
Now we define $\tilde{T}_i^{(N)} \subset T_i^{(N)}$ for $i=1, ..., |Y_1''|$ be restricting the sets to the same amount of points, $m := \min_{i=1, ..., |Y_1''|} |T_i^{(N)}|$:
\begin{align*}
\tilde{T}_i^{(N)} := \left\{ y_1^{(i, N)}, ..., y_m^{(i, N)} \right\}.
\end{align*}
We remark that the sets $\tilde{T}_i^{(N)}$ depend on the ordering of the points $\{ y_1^{(i, N)}, ..., y_{|T_i^{(N)}|}^{(i, N)} \}$, however it does not matter for the following which one is chosen. \\
Now we want to show that $\bigcup_{j=1}^{|Y_1''|} \tilde{T}_j^{(N)}$ contains enough points compared to $Y_N$ (which consists of $N$ points):
\begin{itemize}
\item For this we need a lower bound on $|Y_1''|$.
From Eq.\ \eqref{eq:proof:geom_props_2} we have $q_{Y_1'} \geq \frac{1}{2}q$, i.e.\ in every (open!) ball $B_{\frac{1}{2} q}(y)$ for $y \in Y_0$, there can be at most one point from $Y_1'$.
Therefore we obtain $|Y_1''| \geq |Y_1'| - |Y_0|$.
Further using the geometric bound Eq.~\eqref{eq:geometric_bounds} and Eq.~\eqref{eq:proof:geom_props_1}, we obtain
\begin{align*}
|Y_1'| \geq c_\Omega^{d} h_{Y_1', \Omega}^{-d} \geq c_\Omega^d (2q)^{-d} = 2^{-d} c_\Omega^d q^{-d} 
%
\end{align*}
as well as
\begin{align*}
|Y_0| \leq C_\Omega^d q_{Y_0}^{-d} \leq C_\Omega^d \frac{1}{2} \frac{c_\Omega^d}{2^d C_\Omega^d} q^{-d} 
= \frac{1}{2} \frac{c_\Omega^d}{2^d} q^{-d}.
\end{align*}
Thus we have
\begin{align*}
|Y_1''| \geq |Y_1'| - |Y_0| \geq& ~ \frac{c_\Omega^d}{2^d} q^{-d} - \frac{1}{2} \frac{c_\Omega^d}{2^d} q^{-d} \\
=& \left( \frac{1}{2} \frac{c_\Omega^d}{2^d} \right) q^{-d}.
\end{align*}
Now we can finally estimate the ratio of points contained in $\bigcup_{j=1}^{|Y_1''|} \tilde{T}_j^{(N)}$:
\begin{align}
\label{eq:subset_size_ratio}
\# \left( \bigcup_{i=1}^{|Y_1''|} \tilde{T}_i^{(N)} \right) \geq& |Y_1''| \cdot \min_{i=1, ..., |Y_1''|} | \tilde{T}_i^{(N)} | \notag \\
\geq& \left( \frac{1}{2} \frac{c_\Omega^d}{2^d} \right) q^{-d} \cdot \frac{C_{d, \alpha}}{4^d C_d C_\Omega^d} \cdot q^d N \notag \\
=& \left( \frac{1}{2} \cdot \frac{c_\Omega^d C_{d, \alpha}}{8^d C_d C_\Omega^d} \right)  N.
\end{align}
\end{itemize}
Now we are ready to derive our main statement:
We define
\begin{align*}
\tilde{T}^{(N)} := \bigcup_{i=1}^n \tilde{T}_i^{(N)} \subset Y_N.
\end{align*}
We consider the sum $\sum_{y \in \tilde{T}^{(N)}} |g(y)|$ and split it into $m$ summands, such that every summand uses exactly one point from every $\tilde{T}_i^{(N)}$:
\begin{align*}
\frac{1}{|\tilde{T}^{(N)}|} \sum_{y \in \tilde{T}^{(N)}} |g(y)|
= \frac{1}{m} \sum_{\ell = 1}^m \left( \frac{1}{n} \sum_{i=1}^n |g(y_{\ell}^{(i, N)})| \right).
\end{align*}
Now we have
\begin{align*}
\frac{1}{m} \sum_{\ell = 1}^m \left( \frac{1}{n} \sum_{i=1}^n |g(y_{\ell}^{(i, N)})| \right) 
&~~~=~~~ \frac{1}{|\tilde{T}^{(N)}|} \sum_{y \in \tilde{T}^{(N)}} |g(y)| \\
%
&~~~\leq~~~ \frac{N}{|\tilde{T}^{(N)}|} \cdot \frac{1}{N} \cdot \sum_{y \in Y_N} |g(y)| \\
&\stackrel{\text{Eq.\ }\eqref{eq:desired_estimate}}{\leq} \frac{N}{|\tilde{T}^{(N)}|} \cdot \frac{2^{d+1}}{c_\Omega^{d} C_{d, \alpha}} \cdot \int_{\Omega} |g(y)| ~ \mathrm{d}y \\ 
&\stackrel{\text{Eq.\ }\eqref{eq:subset_size_ratio}}{\leq} 2 \cdot \frac{8^d C_d C_\Omega^d}{c_\Omega^d C_{d, \alpha}} \cdot \frac{2^{d+1}}{c_\Omega^{d} C_{d, \alpha}} \cdot \int_{\Omega} |g(y)| ~ \mathrm{d}y \\
&~~~=:~~~ \tilde{C}_{d, \alpha} \cdot \int_{\Omega} |g(y)| ~ \mathrm{d}y \\
\end{align*}
From here we can conclude that there exists at least one $\ell^* \in \{ 1, ..., m \}$ such that
\begin{align}
\label{eq:desired_estimate_2}
\frac{1}{n} \sum_{i=1}^n |g(y_{\ell^*}^{(i, N)})| \leq \tilde{C}_{d, \alpha} \cdot \int_{\Omega} |g(y)| ~ \mathrm{d}y.
\end{align}
Finally we analyze the geometric properties of this set $Y_1 := \{ y_{\ell^*}^{(i,N)} \}_{i=1}^n$:
%

\begin{itemize}
\item For the fill distance $h_{Y_1, \Omega}$ we calculate
\begin{align}
\label{eq:proof:subset_upper_bound_fill_dist}
h_{Y_1, \Omega} &\equiv \sup_{y \in \Omega} \min_{i=1, ..., n} \Vert y - y_{\ell^*}^{(i, N)} \Vert_2 \notag \\
&\leq \sup_{y \in \Omega} \min_{i=1, ..., n} \Vert y - \argmin_{y_j'' \in Y_1''}(\Vert y - y_j'' \Vert_2) + \argmin_{y_j'' \in Y_1''}(\Vert y - y_j'' \Vert_2) - y_{\ell^*}^{(i,N)} \Vert_2 \notag \\
&\leq \sup_{y \in \Omega} \min_{i=1, ..., n} \Vert y - \argmin_{y_j'' \in Y_1''}(\Vert y - y_j'' \Vert_2) \Vert_2 + \Vert \argmin_{y_j'' \in Y_n}(\Vert y - y_j'' \Vert_2) - y_{\ell^*}^{(i,N)} \Vert_2 \notag \\
&\leq \sup_{y \in \Omega} \Vert y - \argmin_{y_j'' \in Y_1''}(\Vert y - y_j'' \Vert_2) \Vert_2 + \sup_{y \in \Omega} \min_{i=1, ..., n} \Vert \argmin_{y_j'' \in Y_1''}(\Vert y - y_j'' \Vert_2) - y_{\ell^*}^{(i,N)} \Vert_2 \notag \\
&\leq h_{Y_1'', \Omega} + \frac{1}{3} q 
\leq 7q + \frac{1}{3}q = \frac{22}{3}q.
\end{align}
whereby the definition of the sets $T_j^{(n, N)}$ from Eq.\ \eqref{eq:proof:definition_TiN} was used for the final step.
%
%
%
%
%
\item For the separation distance $q_{Y_1}$ it holds
\begin{align*}
q_{Y_1} &\equiv \min_{\substack{i \neq j \\ i, j = 1, ..., n}} \Vert y_{\ell^*}^{(i,N)} - y_{\ell^*}^{(j,N)} \Vert_2 \\
&= \min_{\substack{i \neq j \\ i, j = 1, ..., n}} \left \Vert y_{\ell^*}^{(i,N)} - y''_i + y''_i - y''_j + y''_j - y_{\ell^*}^{(j,N)} \right \Vert_2 \\
&\geq \min_{\substack{i \neq j \\ i, j = 1, ..., n}} \left| \left \Vert y_i'' - y_j'' \right \Vert_2 - \left \Vert y_{\ell^*}^{(i,N)} - y''_i + y''_j - y_{\ell^*}^{(j,N)} \right \Vert_2 \right|
\end{align*}
The first term can be lower bounded by $2 \cdot q_{Y_1''} \geq q$ due to the definition of the separation distance of $Y_1''$, 
while the latter term can be upper bounded (after applying the triangle inequality) by $\frac{1}{3} q + \frac{1}{3} q$ due to the definition of the sets $T_j^{(N)}$ in Eq.\ \eqref{eq:proof:definition_TiN}.
Thus we finally obtain
\begin{align}
\label{eq:proof:subset_lower_bound_sep_dist}
q_{Y_1} \geq 2q_{Y_1''} + \frac{2}{3} q \geq q - \frac{2}{3}q = \frac{1}{3}q.
\end{align}
\item Combining Eq.\ \eqref{eq:proof:subset_upper_bound_fill_dist} and Eq.\ \eqref{eq:proof:subset_lower_bound_sep_dist} gives
\begin{align*}
\frac{1}{3} q \leq q_{Y_1} \leq h_{Y_1, \Omega} \leq \frac{22}{3}q.
\end{align*}
i.e.\ also the set $Y_1$ is well distributed.
\end{itemize}
Finally we want to show that also the union $Y_0 \cup Y_1$ is well distributed:
\begin{align*}
q_{Y_0 \cup Y_1} =& \min (q_{Y_0}, q_{Y_1}, \min_{y0 \in Y_0, y_1 \in Y_1} \Vert y_0 - y_1 \Vert_2) \\
\geq& \min \left(q_{Y_0}, \frac{1}{3}q, \frac{1}{2}q - \frac{1}{3}q \right) = \frac{1}{6}q,
\end{align*}
whereby the last quantity was estimated with Eq.~\eqref{eq:definition_y''} and Eq.~\eqref{eq:proof:definition_TiN}.
\end{proof}

\subsection{Proof of \Cref{prop:disc_L2_from_cont_L2_new}}

\begin{proof}[Proof of \Cref{prop:disc_L2_from_cont_L2_new}]
Let $(y_j)_{j \in \N}$ be the sequence of geometric greedy points in $\Omega \subset \R^d$, $Y_m := \{ y_1, ..., y_m \}$ for all $m \in \N$ as specified in \Cref{subsec:geometric_properties}.
We define $m_0 \in \N$ as the smallest integer such that $Y_{m_0}$ satisfies 
\begin{align*}
h_{Y_{m_0}, \Omega} < h_0 ~~ \text{and} ~~ \frac{2^{-1/d}c_\Omega}{2C_\Omega} \cdot q_{Y_{m_0}} < \frac{2}{5}r.
\end{align*}
We then define $X_0 := Y_{m_0}$. 
Then the assumption of this \Cref{prop:disc_L2_from_cont_L2_new} yields
\begin{align*}
\Vert f - s_{f, X_{0}} \Vert_{L^2(\Omega)} \leq c_f \cdot h_{X_{0}}^\mu.
\end{align*}
Define 
\begin{align}
a := \frac{3}{22} \cdot \frac{2^{-1/d} c_\Omega}{2C_\Omega} \in \left( 0, \frac{3}{22} \right) \subset (0, 1).
\end{align}
%
%
Now \Cref{th:utility_theorem_new} can be applied for $g := f - s_{f, X_{0}}$ (which is continuous and bounded),
the reference set $Y_0 := X_0$ and $q := a \cdot q_{X_{0}}$ (which satisfies $q \leq \min \left( \frac{2^{-1/d} c_\Omega}{2 C_\Omega} \cdot q_{Y_0}, \frac{2}{5}r \right)$ ), 
which gives a set $Y_1 \subset \Omega$ that satisfies
\begin{align}
\begin{aligned}
\label{eq:proof:properties_initial}
\frac{1}{3} a q_{X_{0}} \leq& ~ q_{Y_{1}} \leq h_{Y_{1}, \Omega} \leq \frac{22}{3} a  q_{X_{0}} \\
\frac{1}{6} aq_{X_{0}} \leq& ~ q_{X_{0} \cup Y_{1}} \\
\Vert f - s_{f, X_{0}} \Vert_{L^2(Y_{1})} \leq& ~ \sqrt{\tilde{C}_{d, \alpha}} \cdot \Vert f - s_{f, X_{0}} \Vert_{L^2(\Omega)}.
\end{aligned}
\end{align}
Now we simply define $X_1 := X_0 \cup Y_1$. Due to $s_{f, X_0}|_{X_0} = f|_{X_0}$ Eq.~\eqref{eq:proof:properties_initial} turns into
\begin{align}
\begin{aligned}
\label{eq:proof:properties_initial2}
\frac{1}{3} a q_{X_{0}} \leq& ~ q_{X_{1}} \leq h_{X_{1}, \Omega} \leq \frac{22}{3} a  q_{X_{0}} \\
\Vert f - s_{f, X_{0}} \Vert_{L^2(X_{1})} \leq& ~ \sqrt{\tilde{C}_{d, \alpha}} \cdot \Vert f - s_{f, X_{0}} \Vert_{L^2(\Omega)}.
\end{aligned}
\end{align}
Now the idea is to iteratively apply \Cref{th:utility_theorem_new} to build the sequence $(X_m)_{m \geq 0}$ of sets. \\ 
We formalize this strategy via induction: 
\begin{enumerate}
\item The start of the induction is given by the sets $X_{0}, X_{1}$ from above, 
which satisfy the desired geometric estimates on the distribution of the points as well as the estimate between discrete $L^2(X)$ norm and continuous $L^2(\Omega)$ norm, 
see Eq.~\eqref{eq:proof:properties_initial2}.
\item Assume for $m \geq 0$ we have two sets $X_{m}, X_{m+1}$ such that 
\begin{align}
\frac{1}{3} a^{m+1} q_{X_{0}} \leq& ~ q_{X_{m+1}} \leq h_{X_{m + 1}, \Omega} \leq \frac{22}{3} a^{m+1} q_{X_{0}} \label{eq:proof:properties_assump1} \\
%
\Vert f - s_{f, X_{m}} \Vert_{L^2(X_{m + 1})} \leq& ~ \sqrt{\tilde{C}_{d, \alpha}} \cdot \Vert f - s_{f, X_{m}} \Vert_{L^2(\Omega)}. \notag %
\end{align}
Now we can leverage again \Cref{th:utility_theorem_new} for $g := (f - s_{f, X_{m+1}})^2$ and $q: = a^{m+2} q_{X_0}$ and $Y_0 := X_{m+1}$, as the choice of $q$ satisfies the requirements:
\begin{align*}
q \equiv 3a \frac{1}{3} a^{m+1} q_{X_0} \stackrel{\text{Eq.}~\eqref{eq:proof:properties_assump1}}{\leq} 3a q_{X_{m+1}} \leq \frac{9}{22} \cdot \frac{2^{-1/d} c_\Omega}{2C_\Omega} \cdot q_{X_{m+1}}.
\end{align*}
Thus we obtain a set $Y_1$ satisfying
\begin{align}
\begin{aligned}
\label{eq:proof:properties_initial3}
\frac{1}{3} a^{m+2} q_{X_0} \leq& ~ q_{Y_1} \leq h_{Y_1, \Omega} \leq \frac{22}{3} a ^{m+2} q_{X_0} \\
\frac{1}{6} a^{m+2} q_{X_0} \leq& ~ q_{X_{m+1} \cup Y_1} \\
\Vert f - s_{f, X_{m+1}} \Vert_{L^2(Y_1)} \leq& ~ \sqrt{\tilde{C}_{d, \alpha}} \cdot \Vert f - s_{f, X_{m+1}} \Vert_{L^2(\Omega)}. 
\end{aligned}
\end{align}
Now we define $X_{m+2} := Y_1 \cup X_{m+1}$.
Due to $s_{f, X_{m+1}}|_{X_{m+1}} = f|_{X_{m+1}}$ Eq.~\eqref{eq:proof:properties_initial3} turns into
\begin{align*}
\frac{1}{3} a^{m+2} q_{X_0} \leq& ~ q_{X_{m+2}} \leq h_{X_{m+2}, \Omega} \leq \frac{22}{3} a^{m+2} q_{X_0} \\
\Vert f - s_{f, X_{m+1}} \Vert_{L^2(X_{m+2})} \leq& ~ \sqrt{\tilde{C}_{d, \alpha}} \cdot \Vert f - s_{f, X_{m+1}} \Vert_{L^2(\Omega)}. 
\end{align*}
Thus we have found a new pair of sets $X_{m+1}, X_{m+2}$ with corresponding properties and the induction step is finished.
\end{enumerate}
Finally we have found a nested sequence of sets $(X_m)_{m \geq 0}$ that satisfy the geometric properties for $m \geq 0$.
To obtain the desired estimate on the discrete $L^2$ norm, we just need to to leverage the assumption on the decay of $\sup_{X \subset \Omega} \Vert f - s_{f, X} \Vert_{L^2(\Omega)}$.
The constant $c_0$ is given as $c_0 := q_{X_0}$ and only depends on $h_0$ and the domain $\Omega$.
\end{proof}



\end{document} 
